% =============================================================================
% DSA4262 final report - Group Foursight
%
% SECOND DRAFT (17 Oct 2026).
%
% TO DO before submission (28 Oct):
%   1. Task 2. DONE: SG-NEx methods (Section 3.3), Results (2), Discussion 6.2
%      and the abstract's SG-NEx sentence are written. The "Score cut-offs
%      across cell lines" appendix table was dropped (not used).
%   2. Page limit. Confirmed: the appendix does NOT count. Only the body must
%      fit 5 pages (title page, contents, references and appendix excluded).
%      Once the content is final, optionally change floats from [H] to [htbp].
%   3. Use of AI Methods. DONE: the tools table covers each member's AI tool
%      use; the errors table keeps three cases -- co-modification retraction,
%      the parallel-code bug and the unapproved agent jobs (renumbered 1-3) --
%      and the caption's instruction has been removed.
%   4. Code availability. DONE: both URLs confirmed live and the dashboard
%      content confirmed. Re-confirm the README prediction command if it moves.
%   5. DONE: shipping rank-only (no calibration or F1/threshold table); the
%      related CHECK comments have been removed.
%
% Compile with XeLaTeX or LuaLaTeX (fontspec + Arial):
%     cd report && latexmk -xelatex report.tex
%
% Conventions:
%   - "Cell line 1" = data0 (HCT116, the course's training file).
%     "Cell line 2" = data1 (HEK293T, embryonic kidney; own labels).
%     data2 = synthetic sequences. Cell lines confirmed by course staff (14 Oct).
%   - Every number is held-out (gene-grouped 5-fold split, seed 4262) unless
%     the text says "in-sample".
%   - % SOURCE: comments name the result file each number comes from.
%   - % CHECK: comments flag things to verify before submission.
% =============================================================================
% Must precede \documentclass: makes the Computer Modern maths fonts scalable,
% so maths at the pipeline figure's 6.5/7.5 pt sizes is not substituted.
\RequirePackage{fix-cm}
\documentclass[11pt,a4paper]{article}
% Margin: see the LAYOUT block below for every spacing setting.
\usepackage[a4paper,margin=2cm]{geometry}
\usepackage{fontspec}
\setmainfont{Arial}

\usepackage{graphicx}
\usepackage{booktabs}
\usepackage{tabularx}
\usepackage{array}
\newcolumntype{Y}{>{\raggedright\arraybackslash}X}
\usepackage{amsmath,amssymb,mathtools}
% Blackboard-bold digits (the indicator 1): amssymb's \mathbb has capitals only.
\usepackage{bbm}
\usepackage[round]{natbib}
\usepackage{float}
\usepackage{caption}
\usepackage{subcaption}
\usepackage{tikz}
\usepackage{xcolor}
\usepackage{hyperref}
\usetikzlibrary{shapes.geometric,arrows.meta,positioning,calc,fit,decorations.pathreplacing}

\hypersetup{hidelinks}
\setlength{\emergencystretch}{3em}
\DeclareCaptionFont{captionnine}{\fontsize{9pt}{11pt}\selectfont}
\captionsetup{
    font={captionnine,it},
    labelfont={color=gray},
    textfont={color=gray},
    skip=4pt
}

% --- LAYOUT: every spacing setting in one place ------------------------------
% The brief fixes the font (Arial 11; 9-11 for captions, tables, references)
% but not margins or spacing. Margin is set in the geometry line above (2 cm;
% the template had 2.54 cm). Tune the values here; delete a line to get
% LaTeX's default back.
\usepackage{titlesec}
% \titlespacing*{command}{left indent}{space before}{space after}
\titlespacing*{\section}{0pt}{1.4ex plus .5ex minus .2ex}{0.7ex plus .2ex}
\titlespacing*{\subsection}{0pt}{1.1ex plus .4ex minus .2ex}{0.5ex plus .2ex}
\titlespacing*{\subsubsection}{0pt}{0.9ex plus .3ex minus .2ex}{0.4ex plus .1ex}
\titlespacing*{\paragraph}{0pt}{0.7ex plus .3ex minus .2ex}{0.8em}
% Space around figures and tables: [H] floats use \intextsep.
\setlength{\intextsep}{6pt plus 2pt minus 2pt}
\setlength{\floatsep}{6pt plus 2pt minus 2pt}
\setlength{\textfloatsep}{8pt plus 2pt minus 2pt}
% Lists: tight items, no extra space above and below.
\usepackage{enumitem}
\setlist{itemsep=1pt, parsep=0pt, topsep=2pt, partopsep=0pt}
% Displayed equations (set at \begin{document}, where LaTeX resets them).
\AtBeginDocument{%
    \setlength{\abovedisplayskip}{4pt plus 2pt minus 2pt}%
    \setlength{\belowdisplayskip}{4pt plus 2pt minus 2pt}%
    \setlength{\abovedisplayshortskip}{2pt plus 1pt}%
    \setlength{\belowdisplayshortskip}{2pt plus 1pt}}
% References: 9 pt (allowed by the brief), little space between entries.
\setlength{\bibsep}{1pt plus 0.5pt}
% -----------------------------------------------------------------------------

% Every figure lives in figures/ next to this file (report/figures/ in the
% repo). On Overleaf, upload that folder as "figures".
% \reportfig draws a labelled box instead of failing when an image is missing,
% so a partial upload still compiles.
\newcommand{\reportfig}[2][]{%
    \IfFileExists{figures/#2}%
        {\includegraphics[#1]{figures/#2}}%
        {\fbox{\begin{minipage}[c][3cm][c]{0.9\linewidth}
            \centering\color{gray}\small Missing figure: \texttt{figures/\detokenize{#2}}
        \end{minipage}}}%
}

% Links to the appendix. Extra detail that cannot fit the page limit lives in
% the appendix and is reached by a clickable superscript: write
% \applink{label} in the text and \appmark{label} at the target (a heading or
% caption). The NUMBER is assigned at the target by \appmark, in the order the
% targets appear in the appendix, so the appendix headings read 1, 2, 3, ...
% straight down. \applink looks that number up by label (via the .aux file, so
% it needs the usual second LaTeX pass), which means moving or cutting a
% section never needs renumbering by hand. A second \applink to the same target
% reuses its number.
\newcounter{appmarknum}
\DeclareRobustCommand{\appmark}[1]{%
    \refstepcounter{appmarknum}\label{appmark@#1}%
    \textsuperscript{\arabic{appmarknum}}}
\DeclareRobustCommand{\applink}[1]{%
    \hyperref[#1]{\textsuperscript{\ref*{appmark@#1}}}}

% Model names as they appear in the code and result files.
\newcommand{\mdl}[1]{\texttt{\detokenize{#1}}}
% Shorthand used in tables.
\newcommand{\clone}{CL1}
\newcommand{\cltwo}{CL2}

\title{\textbf{Detection of m6A from nanopore direct RNA sequencing using a site-graph learning framework}}
\author{
    \textbf{DSA4262 Group Project Report -- Team Foursight} \\
    \vspace{0.3cm} \\
    \begin{tabular}{lc}
        \toprule
        \textbf{Student Name} & \textbf{Matriculation Number} \\
        \midrule
        Crescent Liew Qi En & A0277921W \\
        Joshua Jules Gladiston & A0272695N \\
        Robin Ghosh & A0271671A \\
        Siah Jin Thau & A0287931R \\
        \bottomrule
    \end{tabular} \\
    \vspace{0.3cm}
}
\date{}

\begin{document}

\pagenumbering{roman}
\maketitle
% The abstract sits on the title page, which the brief excludes from the 5-page
% limit. Every number is quoted from the body: r(d) Section 3.2.2; PR AUC and
% gains Table 2 / report_significance.csv; m6Anet Table 3.
\begin{abstract}
\noindent Methods that detect m6A from nanopore direct RNA sequencing, m6Anet
among them, score each candidate site independently from its own reads. We
propose a novel model architecture which propagates read-derived evidence 
between candidate sites along a transcript, so that each site's call can be 
corroborated or contradicted by its neighbours. Our final model reaches 0.607
and 0.460 PR AUC on the two provided cell lines (HCT116 and HEK293T), +0.104
and +0.105 over m6Anet (one-sided $p<0.001$). Our paper also shows why a standard graph network fails on this task, traces how we
built a graph architecture suited to it, and investigates how neighbour
evidence can improve performance on unseen cell lines.
% +0.104 / +0.105: results/report_significance.csv (ens_final vs m6anet, both
% arm). The cell line 1 gain understates ours: m6Anet was trained on that line.
Independent GLORI comparisons in three additional cell lines support
SG-NEx candidate prioritisation, while depth-adjusted gene-group analyses
identify shared enrichment patterns across seven cell lines.
\end{abstract}
\vspace{0.2cm}
\setcounter{tocdepth}{1}
\tableofcontents
\clearpage

\pagenumbering{arabic}

% =============================================================================

% =============================================================================
\section{Introduction}

Nanopore direct RNA sequencing enables antibody-free m6A
detection from electrical-current signals. Existing detectors,
including m6Anet \citep{hendra2022}, score candidate sites
independently, limiting the evidence available at low read
depths, as commonly observed in SG-NEx \citep{chen2025}.
Our literature and public-code survey identified no nanopore
RNA detector that propagates read-derived evidence between
candidate sites.\applink{app:related}

Nearby m6A sites are co-modified more often than expected
by chance \citep{m6aiso2025,kang2025}, a pattern also present
in our training labels (Section~\ref{sec:cluster}). We therefore
model candidate sites along each transcript as a graph,
allowing read-derived evidence to propagate between
neighbouring sites. Related approaches for neighbouring DNA methylation sites
are compared in the appendix.\applink{app:deepmod}

Because modification labels differ between our two labelled cell lines, we
investigate which forms of propagation generalise, evaluating held-out genes and testing cross-cell-line
transfer for eligible designs. Task~2 applies the final model to processed SG-NEx samples from seven
cell lines and evaluates predictions against independent GLORI calls.

% =============================================================================
\section{Problem Statement}

We predict m6A modification at candidate DRACH sites from nanopore direct
RNA sequencing. Inputs are m6Anet dataprep outputs: each site's 7-mer and
reads, with 9 current features per read. The model assigns each site
(transcript, position) a score in $[0,1]$, evaluated using ROC AUC and PR
AUC. Because only 4.5\% of sites are modified, \textbf{PR AUC is our headline
metric}.

We investigate whether evidence from nearby sites on the same transcript
improves prediction, and how best to incorporate it. Two requirements guide
the design choices in Section 3.2:
\begin{itemize}
    \item \textbf{Generalisation to unseen cell lines.} We have labels for two
    cell lines, HCT116 and HEK293T, but test data and SG-NEx span others, and
    modification labels vary across cell lines: among shared sites modified in
    at least one of our two labelled cell lines, 64\% are modified in only one. 
    The model must therefore learn transferable patterns rather than 
    cell-line-specific labels.
    \item \textbf{Robustness to read depth.} Our training sites have at least 20 
    reads, whereas SG-NEx sites have a median of 3--5. Because each read 
    represents a separate molecule, lower read depth means less evidence. The 
    model must remain reliable as read depth decreases.
\end{itemize}

% =============================================================================
\section{Methods}
\subsection{Evaluation Design}
\label{sec:data}
We use data0 (HCT116; 121,838 candidate sites, 4.49\% modified) and data1
(HEK293T; 90,810 sites, 7.24\% modified), hereafter cell lines 1 and 2.
Five-fold cross-validation uses one gene partition shared across both cell
lines, preventing overlap of test genes and their transcripts, sites or reads
with training \citep{bbag434}. Within each training fold, 10\% of genes are
reserved for early stopping; label-derived parameters use training genes only.
Models trained on both cell lines are evaluated on held-out genes. Transfer
to an unseen cell line is tested separately by training eligible models on
cell line 1 alone and evaluating on cell line 2, including genes absent from
data0. The final two-head ensemble requires both cell lines for training.

PR AUC is the primary metric, with ROC AUC also reported. Comparisons use
identical held-out sites, with 95\% intervals and one-sided $p$-values from
1,000 whole-gene bootstrap resamples; fold-only comparisons use the
Nadeau--Bengio correction \citep{nadeau2003}.\applink{app:nb}
Robustness is assessed by thinning held-out sites and neighbours to 1, 3 and
10 reads, and by evaluation on two additional HCT116 runs and synthetic RNA
(data2), neither used for training. Protocol and selection details are in
Supplementary Methods.\applink{app:evaluation-details}

\subsection{The Final Model}
\label{sec:model}
The final model combines three architectures, each trained with two random
seeds (Figure~\ref{fig:model_pipeline}).\applink{app:model-overview} A shared DeepSet encoder
\citep{zaheer2017} transforms each read and pools embeddings by their mean
and spread, with log read count included in the site representation. An
auxiliary head predicts modification from the site's own reads.

The two-head H2GCN component \citep{zhu2020} connects candidate sites on the
same transcript within 200\,nt. It keeps the site's embedding, one- and
two-hop neighbour summaries, and a whole-transcript summary separate before
concatenation. Two output heads learn the cell lines' labels separately and
are averaged at inference. The other components use neighbours within
150\,nt: residual + dropout adjusts the site's own-reads score using
distance-weighted neighbour evidence; scalar messages + dropout passes
neighbours' own-reads scores and read counts instead of embedding vectors.
Both hide neighbours randomly during training.\applink{app:designs}

\begin{figure}[H]
    \centering
    \resizebox{\linewidth}{!}{%
% The submitted architecture, drawn left to right (TikZ; uses the document's
% main font). Colours follow the graph: gold = the scored site, green = its
% direct neighbours, blue = sites two hops away.
\definecolor{targetcolor}{HTML}{E8AC59}
\definecolor{onehopcolor}{HTML}{77C29A}
\definecolor{twohopcolor}{HTML}{76A6D2}
\definecolor{ink}{HTML}{333333}
\begin{tikzpicture}[
    font=\fontsize{8pt}{9.6pt}\selectfont, ink, line cap=round,
    stage/.style={font=\fontsize{9pt}{10.8pt}\selectfont\bfseries, anchor=north},
    note/.style={align=center, font=\fontsize{7.5pt}{9pt}\selectfont, inner sep=1pt},
    neuron/.style={circle, draw=black!55, fill=white, minimum size=2.2mm, inner sep=0pt, line width=0.4pt},
    wire/.style={black!30, line width=0.3pt},
    cell/.style={draw=black!50, minimum width=2.4mm, minimum height=2.4mm, inner sep=0pt, line width=0.35pt},
    flow/.style={-{Latex[length=2.2mm,width=1.8mm]}, line width=0.9pt, black!55},
    thin/.style={-{Latex[length=1.6mm,width=1.4mm]}, line width=0.6pt, black!50},
    edge/.style={draw=black!48, line width=0.6pt},
    target/.style={circle, draw=black!65, fill=targetcolor, line width=0.6pt, minimum size=3.9mm, inner sep=0pt,
                   font=\fontsize{6.5pt}{7pt}\selectfont},
    onehop/.style={circle, draw=green!40!black, fill=onehopcolor, line width=0.5pt, minimum size=2.9mm, inner sep=0pt},
    twohop/.style={circle, draw=blue!50!black, fill=twohopcolor, line width=0.5pt, minimum size=2.9mm, inner sep=0pt},
    outnode/.style={circle, draw=black!65, fill=white, minimum size=4.6mm, inner sep=0pt, line width=0.6pt},
]

% =============================================================== 1. DeepSet
\node[stage] at (2.75,4.45) {Read encoder (DeepSet)};
\foreach \y/\lab [count=\r] in {3.45/read 1, 2.55/read 2, 1.05/read $n$} {
    \foreach \j in {0,...,3} { \node[cell, fill=black!6] at (0.25+0.24*\j,\y) {}; }
    \node[note, anchor=east] at (0.08,\y) {\lab};
    \foreach \k in {0,1,2} {
        \node[neuron] (p\r a\k) at (1.45,\y-0.24+0.24*\k) {};
        \node[neuron] (p\r b\k) at (1.9,\y-0.24+0.24*\k) {};
    }
    \foreach \k in {0,1,2} { \foreach \m in {0,1,2} { \draw[wire] (p\r a\k) -- (p\r b\m); } }
    \draw[thin] (1.07,\y) -- (1.3,\y);
    \foreach \j in {0,1,2} { \node[cell, fill=targetcolor!25] (z\r\j) at (2.35+0.24*\j,\y) {}; }
    \draw[thin] (2.02,\y) -- (2.2,\y);
}
\node[note] at (0.6,1.82) {$\vdots$};
\node[note] at (2.6,1.82) {$\vdots$};
\node[note] at (1.67,0.38) {shared network $\phi$\\(same weights for every read)};
\node[circle, draw=black!60, fill=white, minimum size=9mm, align=center, inner sep=0pt,
      font=\fontsize{7pt}{8pt}\selectfont] (pool) at (3.55,2.35) {mean\\spread\\$\log n$};
\foreach \r in {1,2,3} { \draw[thin] (z\r2.east) -- (pool); }
\node[note] at (3.6,1.6) {pool};
\foreach \k in {0,1,2,3} { \node[neuron] (ra\k) at (4.4,1.97+0.25*\k) {}; }
\foreach \k in {0,1,2} { \node[neuron] (rb\k) at (4.85,2.1+0.25*\k) {}; }
\foreach \k in {0,...,3} { \foreach \m in {0,1,2} { \draw[wire] (ra\k) -- (rb\m); } }
\draw[thin] (pool) -- (4.25,2.35);
\node[note] at (4.62,1.55) {$\rho$};
\foreach \j in {0,...,3} { \node[cell, fill=targetcolor] (h\j) at (5.45,1.99+0.24*\j) {}; }
\draw[thin] (4.98,2.35) -- (5.3,2.35);
\node[note, anchor=south] at (5.45,2.95) {site vector};
\node[outnode] (aux) at (5.45,0.75) {$\sigma$};
\draw[thin, dashed] (h0.south) -- (aux);
\node[note, anchor=north] at (5.45,0.45) {auxiliary head:\\own-reads score};

% =============================================================== 2. site graph
\node[stage] at (8.55,4.45) {Site graph};
\begin{scope}[shift={(8.6,2.25)}, scale=0.45]
  \filldraw[fill=twohopcolor!10, draw=blue!55!black, dashed, line width=0.7pt] (0,0) circle (3.05);
  \filldraw[fill=onehopcolor!16, draw=green!45!black, dashed, line width=0.7pt] (0,0) circle (1.55);
  \coordinate (v) at ( 0.00, 0.00);
  \coordinate (a) at (-0.80, 0.85);  \coordinate (b) at ( 0.90, 0.90);
  \coordinate (c) at ( 0.35,-1.15);  \coordinate (d) at (-1.20,-0.15);
  \coordinate (p) at (-1.55, 2.15);  \coordinate (q) at ( 2.25, 1.70);
  \coordinate (r) at ( 2.35,-0.70);  \coordinate (s) at ( 1.65,-2.10);
  \coordinate (t) at (-1.30,-2.30);  \coordinate (u) at (-2.55, 0.25);
  \foreach \n in {a,b,c,d} {\draw[edge] (v) -- (\n);}
  \draw[edge] (a) -- (p); \draw[edge] (b) -- (q); \draw[edge] (b) -- (r);
  \draw[edge] (c) -- (s); \draw[edge] (c) -- (t); \draw[edge] (d) -- (u);
  \node[target] (gv) at (v) {$i$};
  \foreach \n in {a,b,c,d} {\node[onehop] at (\n) {};}
  \foreach \n in {p,q,r,s,t,u} {\node[twohop] at (\n) {};}
\end{scope}
\node[note] at (8.85,0.25) {nodes: candidate sites of one transcript;\\edges: sites within 100--200\,nt};
\draw[flow] (6.15,2.35) -- (6.95,2.35);
\node[note] at (6.55,2.8) {every\\site};

% =============================================================== 3. channels
\node[stage] at (12.35,4.45) {Separate channels};
\foreach \y/\col/\lab [count=\c] in {3.35/targetcolor/{the site itself},
                                     2.55/onehopcolor/{neighbours (1 hop)},
                                     1.75/twohopcolor/{neighbours (2 hops)},
                                     0.95/black!12/{whole transcript}} {
    \foreach \j in {0,...,3} { \node[cell, fill=\col] (c\c\j) at (11.15+0.24*\j,\y) {}; }
    \node[note, anchor=west] at (12.05,\y) {\lab};
}
\draw[thin] (8.78,2.42) to[out=55, in=180] (c10.west);
\draw[thin] (9.25,2.62) to[out=15, in=180] (c20.west);
\draw[thin] (9.95,2.0) to[out=-5, in=180] (c30.west);
\draw[thin, dashed] (9.6,1.3) to[out=-30, in=180] (c40.west);
\node[note] at (13.3,0.35) {kept apart, then concatenated};
\draw[decorate, decoration={brace, amplitude=4pt}, black!55] (14.55,3.55) -- (14.55,0.75);

% =============================================================== 4. heads
\node[stage] at (16.55,4.45) {Two heads};
\foreach \k in {0,1,2,3} { \node[neuron] (ma\k) at (15.0,1.8+0.25*\k) {}; }
\foreach \k in {0,1,2} { \node[neuron] (mb\k) at (15.45,1.92+0.25*\k) {}; }
\foreach \k in {0,...,3} { \foreach \m in {0,1,2} { \draw[wire] (ma\k) -- (mb\m); } }
\draw[thin] (14.7,2.15) -- (14.88,2.15);
\node[outnode] (h1) at (16.35,2.95) {$\sigma$};
\node[outnode] (h2) at (16.35,1.35) {$\sigma$};
\draw[thin] (15.55,2.3) -- (h1);
\draw[thin] (15.55,2.0) -- (h2);
\node[note, anchor=south] at (16.35,3.2) {cell line 1\\labels};
\node[note, anchor=north] at (16.35,1.1) {cell line 2\\labels};
\node[circle, draw=black!60, fill=white, minimum size=7mm, inner sep=0pt, align=center,
      font=\fontsize{7pt}{8pt}\selectfont] (avg) at (17.55,2.15) {mean};
\draw[thin] (h1) -- (avg);
\draw[thin] (h2) -- (avg);
\node[note] at (17.55,1.35) {averaged at\\prediction};

% =============================================================== output
\draw[flow] (17.95,2.15) -- (18.5,2.15);
\node[note, anchor=west, align=left] at (18.55,2.15) {score per site\\[2pt]\textit{6 trained networks;}\\\textit{final score = their}\\\textit{average rank}};
\end{tikzpicture}}
    \caption{The submitted architecture, from a site's reads (left) to its
    score (right).}
    \label{fig:model_pipeline}
\end{figure}

\paragraph{Training and prediction.}
Networks use positive-class-weighted binary cross-entropy and AdamW, with
whole-transcript batches and early stopping on validation PR AUC. The six
networks' within-file ranks are averaged into a score in $[0,1]$, representing
relative rank rather than calibrated probability. Optimisation and
implementation details are in Supplementary Methods.\applink{app:training-details}

\paragraph{Comparators and supporting experiments.}
Comparators are logistic regression on per-site feature means, standard
deviations and the 7-mer; quantiles + LightGBM \citep{ke2017}; and pretrained
m6Anet \citep{hendra2022}. Screening, clustering normalisation and ablation
procedures are documented in Supplementary Methods.\applink{app:methods-extra}


\subsection{SG-NEx Application and GLORI Evaluation}
\label{sec:t2methods}
We pooled SG-NEx reads within each of seven cell lines and analysed the
ensemble's mean raw output, rather than its submission rank; neither is a
calibrated probability. Genomic candidates were evaluated against independent
GLORI calls in four lines, with HCT116 training sites excluded. GLORI was
used only for evaluation; unlisted candidates in genes with a matched call
were inferred negatives. Comparisons with motif-only ranking and m6Anet
were stratified by read depth; A549, H9 and HepG2 were unseen by both models.
Mapping and evaluation details are in Supplementary Methods.
\applink{app:glori-protocol}

Gene-group analysis compared each gene's proportion of sites in its depth
band's top 5\%, with an independent depth-adjusted GLORI check. The 54
class comparisons are exploratory, without multiplicity correction.
\applink{app:sgnex-gene-methods}
% SOURCE: report/sgnex_glori_findings.md and its generated analysis summaries.

\section{Results (1): Model Evaluation}
\label{sec:results1}
\subsection{Predictive Performance and Cross-Cell-Line Transfer}
\label{sec:perf}
The final ensemble outperformed both baselines on held-out genes from
both cell lines (paired gene-bootstrap $p<0.001$;
Table~\ref{tab:model_comparison}). Models were trained on both lines;
pretrained m6Anet's HCT116 result is not held out.
\applink{app:methods-extra}

\begin{table}[H]
    \centering\small
    \begin{tabular}{lcccc}
        \toprule
        & \multicolumn{2}{c}{\textbf{Cell line 1 (HCT116)}} & \multicolumn{2}{c}{\textbf{Cell line 2 (HEK293T)}} \\
        \cmidrule(lr){2-3}\cmidrule(lr){4-5}
        \textbf{Model} & \textbf{PR AUC} & \textbf{ROC AUC} & \textbf{PR AUC} & \textbf{ROC AUC} \\
        \midrule
        Random classifier & 0.045 & 0.500 & 0.072 & 0.500 \\
        Simple baseline (logistic regression) & 0.415 & 0.899 & 0.331 & 0.813 \\
        Quantiles + LightGBM & 0.472 & 0.918 & 0.392 & 0.829 \\
        m6Anet (pretrained on HCT116) & 0.503$^\dagger$ & 0.932$^\dagger$ & 0.356 & 0.819 \\
        DeepSet (reads only, Section~\ref{sec:reads}) & 0.541 & 0.936 & 0.429 & 0.840 \\
        \textbf{Final model} & \textbf{0.607} & \textbf{0.958} & \textbf{0.460} & \textbf{0.865} \\
        \bottomrule
    \end{tabular}
    \caption{Held-out performance, gene-grouped 5-fold, trained on both cell
    lines. $^\dagger$Not held out: m6Anet was trained on this cell line's data.}
    \label{tab:model_comparison}
\end{table}

Trained on HCT116 alone, both constrained components exceeded m6Anet on
HEK293T ($p<0.001$), including genes absent from data0
(Table~\ref{tab:m6anet_gain}). The full two-head ensemble requires both
cell lines for training.

\begin{table}[H]
    \centering\small
    \begin{tabular}{lcccc}
        \toprule
        & & & \multicolumn{2}{c}{\textbf{PR AUC gain over m6Anet}} \\
        \cmidrule(lr){4-5}
        \textbf{Model (trained on cell line 1 only)} & \textbf{PR AUC} & \textbf{ROC AUC} & \textbf{all sites} & \textbf{new genes} \\
        \midrule
        m6Anet (pretrained on HCT116) & 0.356 & 0.819 & -- & -- \\
        Quantiles + LightGBM & 0.372 & 0.815 & +0.016 [.007, .025] & +0.022 [.002, .040] \\
        Residual + dropout & 0.429 & 0.843 & +0.073 [.064, .082] & +0.081 [.059, .101] \\
        Scalar messages + dropout & 0.431 & 0.833 & +0.076 [.066, .085] & +0.076 [.053, .099] \\
        \bottomrule
    \end{tabular}
    \caption{Scored on cell line 2, a cell line none of these models was
    trained on. Brackets: 95\% gene-resampled intervals.}
    \label{tab:m6anet_gain}
\end{table}

\subsection{Learned Read Features and Neighbour Corroboration}
\label{sec:dev}
\paragraph{Read representations.}\label{sec:reads}
Site-level supervision produced more useful modification representations
than unsupervised reconstruction. Attention pooling, a Set Transformer
and a larger residual DeepSet failed the predefined improvement criterion.
DeepSet trained on cell line 1 also improved unseen-cell-line PR AUC over
LightGBM (0.395 versus 0.372).
\applink{app:read-findings}

\paragraph{Local clustering.}\label{sec:cluster}
After within-transcript normalisation, sites within 25\,nt
were co-modified approximately 2.2 times more often than
expected in both cell lines. Enrichment weakened with
distance and approached the null at 200--400\,nt.
\applink{app:cluster-findings}

\paragraph{Neighbour evidence.}\label{sec:graph}
GCN averaging diluted the rare modified sites' own evidence, performing
approximately 0.14 and 0.10 PR AUC below LightGBM. Separating own-site and
neighbour representations with H2GCN improved PR AUC by 0.010--0.016 over
DeepSet alone, and an auxiliary own-reads head improved it further.
\textbf{Replacing real neighbours with unrelated sites removed most of the gain,
supporting corroboration rather than simply extra features.}\applink{app:graph-findings}

\subsection{Local Constraints Improve Transfer}
\paragraph{Context and flexibility.}\label{sec:reach}\label{sec:capacity}
Increasing the radius from 50 to 400\,nt raised the cell line 1 gain over
LightGBM from +0.081 to +0.112, whereas unseen-cell-line performance peaked
at 100\,nt (+0.043) and declined at 200--400\,nt (+0.032--0.034).
Likewise, the graph transformer improved on H2GCN within cell line 1
(+0.025, $p<0.001$) but not on the unseen line ($-0.003$, $p=0.74$).
\applink{app:reach-capacity}

\paragraph{Constrained messages.}\label{sec:constrain}
Relative to H2GCN, residual corrections with neighbour
dropout and scalar messages improved unseen-cell-line
PR AUC by 0.017 and 0.020, respectively
(both $p<0.001$). Ablations implicated neighbour dropout as an important
contributor to the residual design's advantage, while
neither the depth gate nor the precise distance-weighting
profile showed a clear additional benefit.
\applink{app:constraint-findings}

\subsection{Robustness Across Read Depths and Datasets}
\label{sec:robust}
Under random thinning of both target and neighbour reads, the final model
retained higher PR AUC than LightGBM at every tested depth
(Table~\ref{tab:depth}). Much of the sparse-read advantage came from the
learned encoder; neighbours helped most when better covered than the target.
These experiments simulate reduced coverage and do not establish accuracy
on naturally low-expression sites.\applink{app:robust-findings}

\begin{table}[H]
    \centering\small
    \begin{tabular}{lcccc}
        \toprule
        \textbf{PR AUC} & \textbf{1 read} & \textbf{3 reads} & \textbf{10 reads} & \textbf{All ($\ge$20)} \\
        \midrule
        Final model, \clone & 0.333 & 0.451 & 0.534 & 0.607 \\
        Quantiles + LightGBM, \clone & 0.140 & 0.221 & 0.368 & 0.472 \\
        Final model, \cltwo & 0.279 & 0.341 & 0.397 & 0.460 \\
        Quantiles + LightGBM, \cltwo & 0.136 & 0.200 & 0.306 & 0.392 \\
        \bottomrule
    \end{tabular}
    \caption{Held-out PR AUC with each site's (and its neighbours') reads
    thinned at random; models trained at full depth on both cell lines.}
    \label{tab:depth}
\end{table}

On two further HCT116 runs, a close ensemble variant using scalar messages
without dropout lost 0.009--0.015 PR AUC relative to data0's run.
On synthetic data2, ensemble scores correlated moderately
with modification fraction (Spearman $\rho=0.41$) and
plateaued above 50\% modification.\applink{app:robust-findings}

\section{Results (2): Application on the SG-NEx Data}
\label{sec:sgnex-results}
\subsection{Independent Validation and m6Anet Comparison}
Against independent GLORI calls, the model exceeded motif-only ranking in
every depth band of the four evaluated lines. In unseen A549, H9 and HepG2,
PR AUC at 1--2 pooled reads was 0.135, 0.041 and 0.022, versus 0.073, 0.021
and 0.012 for motif-only ranking. This supports candidate prioritisation at
naturally low depth, without isolating the contribution of neighbours.
At 20+ reads, the model exceeded m6Anet on H9 and HepG2, with no clear
difference on A549 (Table~\ref{tab:glori-main}). GLORI provides incomplete
labels: unlisted candidates are inferred negatives, not verified unmodified
sites.\applink{app:glori-results}

\begin{table}[H]
\centering\small
\begin{tabular}{lccc}
\toprule
\textbf{Cell line} & \textbf{Our PR AUC} & \textbf{m6Anet PR AUC} & \textbf{Paired difference [95\% interval]} \\
\midrule
A549 & 0.366 & 0.357 & +0.008 [$-0.001$, +0.017] \\
H9 & 0.300 & 0.262 & +0.038 [+0.032, +0.044] \\
HepG2 & 0.381 & 0.278 & +0.103 [+0.087, +0.119] \\
\bottomrule
\end{tabular}
\caption{Comparison against GLORI calls on matched genomic candidates with
20+ pooled reads in three unseen cell lines. Intervals use 200 gene-bootstrap
resamples; differences use unrounded values. Unlisted candidates in
GLORI-covered genes are inferred negatives.}
\label{tab:glori-main}
\end{table}

\subsection{Depth-Adjusted Gene-Group Patterns}
Here, enrichment means a greater proportion of candidate sites ranked in
their depth band's top 5\%, not a higher fraction of modified molecules.
Across seven cell lines, transcription factors, TNF$\alpha$/NF-$\kappa$B
signalling, early oestrogen response and long non-coding RNAs were enriched
for high-scoring sites after depth adjustment. MYC and E2F targets, mTORC1
signalling, mitotic spindle and protein secretion were depleted. GLORI
supported the directions in all four evaluated lines under the stated
consistency criteria. These associations are exploratory: 54 classes were
tested without multiplicity correction, and motif composition was not
separated from gene-group effects.\applink{app:sgnex-gene-methods}

Ribosomal-protein genes had high raw scores but fewer depth-adjusted
high-scoring sites than other genes. GLORI did not consistently support
depletion, illustrating coverage confounding rather than confirmed
hypomodification.
\applink{app:glori-fractions}

\section{Discussion}

\subsection{Interpretation of Model Performance}
\label{sec:disc-model}

\paragraph{Sources of predictive improvement.}
Learned read representations account for much of the
low-depth advantage, while neighbouring sites provide
additional evidence. Successful graph designs preserve
the site's own evidence through separate channels or
residual corrections.

\paragraph{Shared local patterns and discordant site labels.}
Similar local co-modification patterns across cell lines
may explain why constrained neighbour evidence transferred
better than broader context. The observed depletion near
exon junctions is consistent with previous studies
\citep{uzonyi2023,he2023}, although its contribution to
clustering and transfer remains uncertain.

At discordantly labelled sites, our models and raw-feature
analyses found little evidence distinguishing the cell
lines. This could reflect experimental labelling differences
or modification-fraction differences unresolved by the
available features and models.

\subsection{Interpretation of the SG-NEx Analysis}
Independent GLORI agreement supports the use of SG-NEx predictions for
candidate prioritisation beyond motif identity, including at naturally low
read depth. Validation in A549, H9 and HepG2 broadens the original two-line
evaluation to additional cell lines and a different assay. It does not show
that neighbour propagation alone caused the transfer benefit; that question
requires matched graph ablations on these data.

The gene-group results emphasise shared patterns across cell lines rather
than strong cell-line-specific differences. Enrichment of regulatory groups
and depletion of several proliferation and biosynthesis groups are
consistent across model predictions and GLORI, but their mechanisms remain
uncertain. The analyses do not establish transcript decay effects, and
motif composition has not been separated from gene-group membership.
The ribosomal-protein example shows why coverage adjustment is essential
before interpreting raw-score differences as biology.

GLORI is an independent measurement, not complete ground truth. Unlisted
sites may be undetected positives; studies used different cultures, and the
HepG2 and HCT116 data came from engineered control-guide cells. Consequently,
absolute performance and model comparisons should be interpreted with these
labelling and sampling limitations. The score--fraction association is useful
for ranking, but does not support quantitative stoichiometry estimates.

\subsection{Limitations and Future Work}
\label{sec:limits}

\begin{itemize}
    \item \textbf{Incomplete external labels.} GLORI publishes positive
    calls rather than verified negative sites. Restricting evaluation to
    genes with a detected call reduces, but does not remove, uncertainty
    about unlisted sites. Four of seven SG-NEx lines have GLORI data;
    HCT116 is not an unseen cell line.
    \item \textbf{Exploratory gene-group comparisons.} The 54 tested
    classes were not corrected for multiple comparisons. Agreement across
    lines and assays supports descriptive consistency, not causal effects
    or independence from sequence composition.

    \item \textbf{Cell line or labelling.} The two labelled files differ in
    cell line \emph{and} in labelling; part of what we call cell-line-specific
    could be how each was labelled.
    \item \textbf{Low-depth validation is incomplete.} The original
    training labels cover sites with at least 20 reads, so their low-depth
    evaluation uses thinning. GLORI adds evidence at naturally low-depth
    sites, but unlisted candidates are inferred negatives and assay detection
    itself depends on coverage.
    \item \textbf{Modification fraction is not directly supervised.}
    The model is trained on binary site labels rather than the
    fraction of molecules carrying m6A. Its scores should therefore
    not be interpreted as modification fractions. On synthetic
    data2, scores plateaued above 50\% modification, limiting
    quantitative estimation.\applink{app:presence} Reliable
    fraction estimation would require additional supervision
    or validation against known modification fractions.
    \item \textbf{Scores indicate relative priority, not probability.}
    The final score averages each network's within-file ranks.
    Higher scores identify sites with stronger model evidence
    for m6A and can prioritise candidates for further analysis.
    Scores should not be used directly to estimate the number
    of modified sites or compare modification prevalence
    across files.
    \item \textbf{Incomplete neighbourhoods.}
    Missing neighbours may reduce performance. Neighbour dropout addresses
    this risk, but robustness of the final ensemble to partial transcripts
    remains uncertain.\applink{app:partial-transcripts}
    \item \textbf{Uncertainty in model comparisons.}
    Most variants were trained with three seeds, with seed
    variability of 0.002--0.005 PR AUC; differences below
    approximately 0.005 should be interpreted cautiously.
    Removing neighbours from a trained network measures its
    reliance on neighbour information, rather than the model's performance when trained without neighbours.
\end{itemize}

\paragraph{Future work.}
Three directions follow from these limitations.
\emph{Validation across assays and cell lines:} comparing predictions with
additional assays and METTL3-knockout data could help distinguish assay
disagreement from biological variation. Testing local co-modification in
more labelled cell lines would assess whether the clustering pattern is
general. \emph{Extending neighbour evidence:} retaining read identities
would allow corroboration between sites on the same molecule
\citep{kang2025,m6aiso2025}, while graphs based on read-signal similarity
could connect sites across transcripts rather than only nearby positions.
\emph{Additional prediction targets:} quantitative modification-fraction
labels could support stoichiometry estimation, independent validation could
assess probability calibration, and isoform-level analyses could examine
modification differences between transcript versions of the same gene.

% =============================================================================
\clearpage
% Supporting material exceeds 26 pages; alphabetic page counters stop at z.
\pagenumbering{arabic}
\renewcommand{\thepage}{S\arabic{page}}

\phantomsection
\addcontentsline{toc}{section}{References}
% Status of each reference is recorded in report/literature.md.
\begin{thebibliography}{99}
    \footnotesize   % 9 pt in an 11 pt document, the brief's minimum for references
    \raggedright    % long titles and DOIs cannot be justified without gaps

    \bibitem[Angermueller et~al.(2017)]{angermueller2017}
    Angermueller, C., Lee, H.\,J., Reik, W., and Stegle, O. (2017).
    \newblock DeepCpG: accurate prediction of single-cell DNA methylation states using deep learning.
    \newblock \textit{Genome Biology}, 18, 67. doi:10.1186/s13059-017-1189-z.

    \bibitem[Liang et~al.(2026)]{bbag434}
    Liang, H., Hua, H., Li, S., and Chen, S. (2026).
    \newblock A systematic benchmarking framework and dual-view optimization strategy for single-cell DNA methylation imputation.
    \newblock \textit{Briefings in Bioinformatics}, 27(4), bbag434. doi:10.1093/bib/bbag434.

    \bibitem[Caruana(1997)]{caruana1997}
    Caruana, R. (1997).
    \newblock Multitask learning.
    \newblock \textit{Machine Learning}, 28, 41--75.

    \bibitem[Chen et~al.(2025)]{chen2025}
    Chen, Y. et~al. (2025).
    \newblock A systematic benchmark of Nanopore long-read RNA sequencing for transcript-level analysis in human cell lines.
    \newblock \textit{Nature Methods}, 22, 801--812. doi:10.1038/s41592-025-02623-4.

    \bibitem[Corso et~al.(2020)]{corso2020}
    Corso, G., Cavalleri, L., Beaini, D., Li\`o, P., and Veli\v{c}kovi\'c, P. (2020).
    \newblock Principal neighbourhood aggregation for graph nets.
    \newblock \textit{NeurIPS}.

    \bibitem[De~Waele et~al.(2022)]{dewaele2022}
    De~Waele, G., Clauwaert, J., Menschaert, G., and Waegeman, W. (2022).
    \newblock CpG Transformer for imputation of single-cell methylomes.
    \newblock \textit{Bioinformatics}, 38(3), 597--603. doi:10.1093/bioinformatics/btab746.

    \bibitem[Feng et~al.(2020)]{feng2020}
    Feng, W. et~al. (2020).
    \newblock Graph random neural networks for semi-supervised learning on graphs.
    \newblock \textit{NeurIPS}.

    \bibitem[Gulrajani and Lopez-Paz(2021)]{gulrajani2021}
    Gulrajani, I. and Lopez-Paz, D. (2021).
    \newblock In search of lost domain generalization.
    \newblock \textit{ICLR}.

    \bibitem[He et~al.(2023)]{he2023}
    He, P.\,C. et~al. (2023).
    \newblock Exon architecture controls mRNA m6A suppression and stability.
    \newblock \textit{Science}. doi:10.1126/science.abj9090.

    \bibitem[Hendra et~al.(2022)]{hendra2022}
    Hendra, C., Pratanwanich, P.\,N., Wan, Y.\,K. et~al. (2022).
    \newblock Detection of m6A from direct RNA sequencing using a multiple instance learning framework.
    \newblock \textit{Nature Methods}, 19, 1590--1598. doi:10.1038/s41592-022-01666-1.

    \bibitem[Huang et~al.(2021)]{huang2021}
    Huang, Q., He, H., Singh, A., Lim, S.-N., and Benson, A.\,R. (2021).
    \newblock Combining label propagation and simple models out-performs graph neural networks.
    \newblock \textit{ICLR}. arXiv:2010.13993.

    \bibitem[Kang et~al.(2025)]{kang2025}
    Kang, G., Hwang, H., et~al. (2025).
    \newblock Comprehensive discovery of m6A sites in the human transcriptome at single-molecule resolution (DeepRM).
    \newblock \textit{Nature Communications}. doi:10.1038/s41467-025-67417-w.
    % 17 authors (Nature Index); co-first authors Gihyeon Kang and Hyeonseo
    % Hwang, corresponding author Daehyun Baek.

    \bibitem[Ilse et~al.(2018)]{ilse2018}
    Ilse, M., Tomczak, J.\,M., and Welling, M. (2018).
    \newblock Attention-based deep multiple instance learning.
    \newblock \textit{ICML}.

    \bibitem[Ke et~al.(2017)]{ke2017}
    Ke, G. et~al. (2017).
    \newblock LightGBM: a highly efficient gradient boosting decision tree.
    \newblock \textit{NeurIPS}.

    \bibitem[Kipf and Welling(2017)]{kipf2017}
    Kipf, T.\,N. and Welling, M. (2017).
    \newblock Semi-supervised classification with graph convolutional networks.
    \newblock \textit{ICLR}.

    \bibitem[Koh et~al.(2019)]{koh2019}
    Koh, C.\,W.\,Q., Goh, Y.\,T., and Goh, W.\,S.\,S. (2019).
    \newblock Atlas of quantitative single-base-resolution N6-methyl-adenine methylomes.
    \newblock \textit{Nature Communications}, 10, 5636. doi:10.1038/s41467-019-13561-z.

    \bibitem[Lee et~al.(2019)]{lee2019}
    Lee, J., Lee, Y., Kim, J., Kosiorek, A., Choi, S., and Teh, Y.\,W. (2019).
    \newblock Set Transformer: a framework for attention-based permutation-invariant neural networks.
    \newblock \textit{ICML}.

    \bibitem[Leger et~al.(2021)]{leger2021}
    Leger, A. et~al. (2021).
    \newblock RNA modifications detection by comparative Nanopore direct RNA sequencing (Nanocompore).
    \newblock \textit{Nature Communications}. doi:10.1038/s41467-021-27393-3.

    \bibitem[Liu et~al.(2019)]{liu2019deepmod}
    Liu, Q. et~al. (2019).
    \newblock Detection of DNA base modifications by deep recurrent neural network on Oxford Nanopore sequencing data (DeepMod).
    \newblock \textit{Nature Communications}, 10. doi:10.1038/s41467-019-10168-2.

    \bibitem[Lorenz et~al.(2011)]{lorenz2011}
    Lorenz, R. et~al. (2011).
    \newblock ViennaRNA Package 2.0.
    \newblock \textit{Algorithms for Molecular Biology}, 6, 26.

    \bibitem[Luo et~al.(2023)]{luo2023}
    Luo, Z. et~al. (2023).
    \newblock Exon-intron boundary inhibits m6A deposition, enabling m6A distribution hallmark, longer mRNA half-life and flexible protein coding.
    \newblock \textit{Nature Communications}, 14, 4172. doi:10.1038/s41467-023-39897-1.

    \bibitem[m6Aiso(2025)]{m6aiso2025}
    Single-molecule m6A detection empowered by endogenous labeling unveils complexities across RNA isoforms (m6Aiso) (2025).
    \newblock \textit{Molecular Cell}. doi:10.1016/j.molcel.2025.01.014.
    % First author NOT confirmed (an AI summary said "Guo et al."): cite by title and DOI.

    \bibitem[Nadeau and Bengio(2003)]{nadeau2003}
    Nadeau, C. and Bengio, Y. (2003).
    \newblock Inference for the generalization error.
    \newblock \textit{Machine Learning}, 52, 239--281.

    \bibitem[Dong et~al.(2026)]{orca2026}
    Dong, H. et~al. (2026).
    \newblock Comprehensive mapping of RNA modification dynamics and crosstalk via deep learning and nanopore direct RNA-sequencing (ORCA).
    \newblock \textit{Nature Communications}. doi:10.1038/s41467-026-68419-y.

    \bibitem[Pratanwanich et~al.(2021)]{pratanwanich2021}
    Pratanwanich, P.\,N. et~al. (2021).
    \newblock Identification of differential RNA modifications from nanopore direct RNA sequencing with xPore.
    \newblock \textit{Nature Biotechnology}. doi:10.1038/s41587-021-00949-w.

    \bibitem[Rampášek et~al.(2022)]{rampasek2022}
    Rampášek, L. et~al. (2022).
    \newblock Recipe for a general, powerful, scalable graph transformer.
    \newblock \textit{NeurIPS}.

    \bibitem[Rong et~al.(2020)]{rong2020}
    Rong, Y., Huang, W., Xu, T., and Huang, J. (2020).
    \newblock DropEdge: towards deep graph convolutional networks on node classification.
    \newblock \textit{ICLR}.

    \bibitem[Saerens et~al.(2002)]{saerens2002}
    Saerens, M., Latinne, P., and Decaestecker, C. (2002).
    \newblock Adjusting the outputs of a classifier to new a priori probabilities: a simple procedure.
    \newblock \textit{Neural Computation}, 14, 21--41.

    \bibitem[Sagawa et~al.(2020)]{sagawa2020}
    Sagawa, S., Koh, P.\,W., Hashimoto, T.\,B., and Liang, P. (2020).
    \newblock Distributionally robust neural networks for group shifts.
    \newblock \textit{ICLR}.

    \bibitem[Sch\"utt et~al.(2017)]{schutt2017}
    Sch\"utt, K.\,T. et~al. (2017).
    \newblock SchNet: a continuous-filter convolutional neural network for modeling quantum interactions.
    \newblock \textit{NeurIPS}.

    \bibitem[Uzonyi et~al.(2023)]{uzonyi2023}
    Uzonyi, A. et~al. (2023).
    \newblock Exclusion of m6A from splice-site proximal regions by the exon junction complex dictates m6A topologies and mRNA stability.
    \newblock \textit{Molecular Cell}, 83(2), 237--251. doi:10.1016/j.molcel.2022.12.026.

    \bibitem[Veličković et~al.(2018)]{velickovic2018}
    Veličković, P. et~al. (2018).
    \newblock Graph attention networks.
    \newblock \textit{ICLR}.

    \bibitem[Xu et~al.(2018)]{xu2018}
    Xu, K. et~al. (2018).
    \newblock Representation learning on graphs with jumping knowledge networks.
    \newblock \textit{ICML}.

    \bibitem[Zaheer et~al.(2017)]{zaheer2017}
    Zaheer, M. et~al. (2017).
    \newblock Deep sets.
    \newblock \textit{NeurIPS}.

    \bibitem[Zhang et~al.(2020)]{zhang2020linpart}
    Zhang, H., Zhang, L., Mathews, D.\,H., and Huang, L. (2020).
    \newblock LinearPartition: linear-time approximation of RNA folding partition function and base-pairing probabilities.
    \newblock \textit{Bioinformatics}, 36, i258--i267.

    \bibitem[Zhu et~al.(2020)]{zhu2020}
    Zhu, J. et~al. (2020).
    \newblock Beyond homophily in graph neural networks: current limitations and effective designs.
    \newblock \textit{NeurIPS}.

    % --- Cited in the appendix only ---
    \bibitem[Gui et~al.(2022)]{gui2022}
    Gui, S., Li, X., Wang, L., and Ji, S. (2022).
    \newblock GOOD: a graph out-of-distribution benchmark.
    \newblock \textit{NeurIPS Datasets and Benchmarks}.

    \bibitem[Niv and Rabin(2025)]{niv2025}
    Niv and Rabin (2025).
    \newblock Exploring graph-transformer out-of-distribution generalization abilities.
    \newblock arXiv:2506.20575.
\end{thebibliography}

\phantomsection
\section*{Code Availability}
\addcontentsline{toc}{section}{Code Availability}

All code, the trained model (\texttt{models/final/}), a small sample dataset
(\texttt{data/sample/}) and installation and prediction instructions are in
the team's public repository. Prediction needs only \texttt{pip install -e .}
and no deep-learning framework; the repository's README documents the
command. Every number in this report names its source script and result file
in \texttt{report/findings.md}.

\noindent\textbf{Code repository:} \url{https://github.com/AsJayTee/DSA4262-Foursight}

\noindent\textbf{Interactive data visualisation platform:}\\ \url{https://asjaytee.github.io/DSA4262-Foursight/}
\phantomsection
\section*{Use of AI Methods}
\addcontentsline{toc}{section}{Use of AI Methods}

\begin{table}[H]
    \centering\small
    \begin{tabularx}{\linewidth}{>{\raggedright\arraybackslash}p{2.3cm}YY}
        \toprule
        \textbf{AI Tool used} & \textbf{Prompt and output} & \textbf{How the output was used in the assignment} \\
        \midrule
        Claude Code (Anthropic) & Asked to implement and run experiments (feature sets, graph networks, evaluation harness, analyses) and to record results; produced code, result files and written summaries. & Code reviewed and tested (\texttt{make smoke}, unit tests pinning numpy inference to torch); every quoted number regenerated from a committed result file; design decisions recorded in \texttt{docs/decisions/}. \\
        ChatGPT deep research & Literature report 1 (10~Oct): nanopore m6A methods and prior art for neighbour-site models, with a log of 31 search strings. & Used to find candidate citations only; each claim checked against the paper before use; several errors found (below). \\
        Claude (deep research) & Literature report 2 (10~Oct): the same prior-art questions, with quoted passages, 26 logged searches and its own list of what it could not verify. & As for report 1: candidates only, every cited claim checked against the paper; its unverified items (e.g.\ NanoFM) not cited. \\
        ChatGPT & Given the report draft (16~Oct) and asked to rephrase it for clarity without changing numbers, LaTeX commands, source comments or hedges; returned suggested rewrites of eight passages and a list of editorial issues. & Each rewrite reviewed against the draft before use: most adopted; two reverted where meaning changed (a weakened novelty claim, a dropped definition of homophily); one flagged inconsistency (a run-to-run cost described as ``about the seed spread'') confirmed and corrected. \\        \bottomrule
    \end{tabularx}
    \caption{AI tools used.}
    \label{tab:ai_usage}
\end{table}

\subsection*{AI Output Errors, Misleading Claims, and Verification}

\begin{table}[H]
    \centering\small
    \begin{tabularx}{\linewidth}{>{\raggedright\arraybackslash}p{1.1cm}YYY}
        \toprule
        \textbf{Case} & \textbf{Incorrect or unverified output} & \textbf{How it was detected} & \textbf{Corrective action} \\
        \midrule
        1 & Reported that co-modification ``fades twice as fast in cell line 2'' (decay scale 59 vs 113\,nt), from point estimates. & Robustness check: transcript-bootstrap intervals overlap (38--123 vs 82--144\,nt); on shared sites the two labellings give 101 vs 103\,nt. & Claim retracted; reported as $\approx$100\,nt in both lines. \\
        2 & AI-written parallel code for the folding test silently ran RNAplfold when LinearPartition was requested (Python 3.14's process start method did not pass the setting to workers). & The ``LinearPartition'' results were identical to RNAplfold's. & Fixed with a worker initialiser; workers now report the predictor they used; rerun. \\
        3 & (Agent action) Launched compute jobs the user had not approved. & The user noticed the extra jobs. & Jobs killed; an explicit approval rule adopted for every compute job. \\
        \bottomrule
    \end{tabularx}
    \caption{Cases where AI output was wrong, misleading or unverified.}
    \label{tab:ai_issues}
\end{table}

% =============================================================================
\clearpage
\phantomsection
\section*{Appendix}
\addcontentsline{toc}{section}{Appendix}
% Lettered so that \ref to the appendix reads 'Appendix A', not the last
% numbered section.
\setcounter{subsection}{0}
\renewcommand{\thesubsection}{\Alph{subsection}}

\setcounter{table}{0}
\renewcommand{\thetable}{A\arabic{table}}
\setcounter{figure}{0}
\renewcommand{\thefigure}{A\arabic{figure}}

% Ordered so the headings read 1, 2, 3, ... straight down: \appmark assigns
% each number at the heading, in the order the headings appear here (see the
% macro definition near the top of the file), and \applink in the body prints
% that number. The sequence roughly follows the order the body first refers to
% each section -- prior art and statistics, then the model and its methods,
% then the Results (1) development findings, then the SG-NEx / GLORI material.
% Only the 20 sections the body actually points to carry a superscript number;
% their detail tables, figures and short derivations live inside them as
% \paragraph material and are reached by ordinary \ref (``Table~A..'',
% ``Figure~A..''), not by a superscript. So every number in the appendix has a
% matching superscript somewhere in the body. To move a section, just move its
% block: the numbers are regenerated on the next compile, so nothing needs
% renumbering by hand.

\phantomsection
\subsection*{Related work and prior art\texorpdfstring{\appmark{app:related}}{}}
\label{app:related}
\paragraph{How we searched.} Between 8 and 10 October 2026 we looked for any
method that shares evidence between candidate modification sites.
\begin{itemize}
    \item \textbf{Sources.} Our own web searches, plus two AI deep-research
    reports run independently on 10~October. Report 1 (ChatGPT) searched
    Google and Bing, PubMed and PubMed Central, arXiv and GitHub, logging 31
    search strings. Report 2 (Claude) ran 26 searches through a general web
    index covering PubMed Central, bioRxiv, arXiv, journal sites and GitHub
    (not native queries to each database).
    \item \textbf{What the searches targeted.} (i) Nanopore RNA modification
    callers, by name and by design: m6Anet, m6ATM, DeepRM, CHEUI, xPore,
    Nanocompore, EpiNano, ELIGOS, Tombo, MINES, SingleMod, TandemMod and
    m6Aiso, plus MultiNano, RNANO, ORCA and Yanocomp in report 2, and
    direct queries for graph neural networks over neighbouring sites in
    nanopore direct RNA data. (ii) The same idea in other data: neighbouring
    CpGs in DNA methylation (nanopore and PacBio callers, single-cell
    imputation, smoothing), hidden-Markov peak callers for MeRIP-seq, and
    sequence- or structure-only m6A predictors. (iii) Single-molecule
    co-occurrence of m6A sites. (iv) Cross-cell-line and cross-assay
    concordance of m6A labels. (v) Graph models under distribution shift,
    and neighbour dropout.
    \item \textbf{Searches that found nothing.} The direct graph queries, in
    both reports, returned no nanopore RNA method; neither report found a
    method that uses single-molecule co-occurrence as a prediction prior in
    RNA. Two independent reports reaching the same verdict is part of our
    evidence.
    \item \textbf{Verification.} The reports were used only to find
    candidates. For every method we cite, we read the paper (full text where
    accessible, otherwise the abstract) to see how it scores a site; the
    reports' errors that this caught are listed in the AI-use section.
\end{itemize}
Every nanopore RNA method found (Table~\ref{tab:related}) scores each site
independently; where co-modification between sites is studied, it is
analysed after prediction, not used to make it. The closest precedents are
for DNA methylation. The searches did not cover conference proceedings
(RECOMB, ISMB), machine-learning workshop papers or theses, so the claim
holds for published papers and public code reachable through these sources
as of 10~October 2026.
% SOURCE: report/literature.md, "Search record" and the full logs of report 1
% (ChatGPT, 31 strings) and report 2 (Claude, 18 searches + 8 by a subagent).

\begin{table}[H]
    \centering\small
    \begin{tabularx}{\linewidth}{>{\raggedright\arraybackslash}p{4.6cm}Y}
        \toprule
        \textbf{Method} & \textbf{What it does} \\
        \midrule
        m6Anet \citep{hendra2022} & Multiple-instance learning over a site's reads; each site scored independently \\
        m6ATM (Brief Bioinform 2024) & Deep network over reads + MIL, per site; RNA structure not used \\
        DeepRM \citep{kang2025} & Transformer over one molecule's signal and a 21-nt context; reports co-occurring m6A on single molecules (post hoc) \\
        Xron (Genome Res 2024) & Methylation-aware basecaller; context within one read \\
        CHEUI, TandemMod, SingleMod, ORCA \citep{orca2026}, MINES, MultiNano, RNANO, Nanocompore \citep{leger2021} & Per-site models; co-occurrence, where studied, analysed after prediction (reported, not individually re-checked) \\
        m6A-IIN (Commun Biol 2025) & Sequence-only: 41-nt fragment + RNAfold structure; no nanopore signal \\
        DeepMod \citep{liu2019deepmod} & Nanopore DNA 5mC; optional second stage revises a CpG's methylation fraction from a histogram of neighbouring CpGs' fractions --- closest precedent \\
        DeepCpG \citep{angermueller2017}, CpG Transformer \citep{dewaele2022}, GraphCpG & Impute single-cell CpG states from neighbouring CpG states \\
        \bottomrule
    \end{tabularx}
    \caption{No nanopore RNA modification method we found shares evidence
    between candidate sites on a transcript; ``graph'' models build graphs
    over nucleotides within one molecule. Searches as of 10~Oct~2026.}
    \label{tab:related}
\end{table}
% SOURCE: report/literature.md (method list, verification status, search record).

\phantomsection
\subsection*{DeepMod vs our model\texorpdfstring{\appmark{app:deepmod}}{}}
\label{app:deepmod}
\textit{Note on how this comparison arose.} Our model was not derived from
DeepMod. We arrived at it from our own data: the co-modification pattern in
the training labels (Section~\ref{sec:cluster}), the corroboration argument
that follows from it (Section~\ref{sec:graph}), and the experiments of
Section~\ref{sec:dev}. We found DeepMod only afterwards, while searching the
literature and public code for any prior model that shares evidence between
nearby sites, to test our claim that no such model exists for nanopore direct
RNA sequencing. It is the closest precedent we found, so we compare the two
here.


\begin{table}[H]
    \centering\small
    \begin{tabularx}{\linewidth}{lYY}
        \toprule
        & \textbf{DeepMod cluster stage (DNA 5mC)} & \textbf{Ours (RNA m6A)} \\
        \midrule
        Neighbourhood & CpGs within $\pm$25\,bp, both strands & Candidate sites on the transcript within 50--200\,nt \\
        Neighbour sends & Final fraction of reads called methylated & Learned read-derived vector, or own-reads score + log read count \\
        Distance & Ignored inside the window & $e^{-d/\lambda}$, $\lambda\approx100$\,nt \\
        Aggregation & Fixed normalised histogram & Learned; weighted mean (+ max within 75\,nt) \\
        Evidence strength & Lost (1/1 and 50/50 reads both give 1.0) & Kept: read counts in messages and the gate \\
        Combination & Network on 14 numbers & Residual: reduces exactly to the own score with no neighbours \\
        Training & Two-stage, on fixed stage-1 outputs & End to end, with neighbour dropout \\
        Setting & Dense, strongly correlated CpGs & Sparse, moderately clustered m6A (2.2$\times$ within 25\,nt); label shift between cell lines \\
        \bottomrule
    \end{tabularx}
    \caption{Both fit $z_i=F(\text{own}_i,\mathrm{AGG}_{j}\,\psi(\text{evidence}_j,d_{ij}))$.
    DeepMod's LSTM is its per-read signal model, the analogue of our read
    encoder; its neighbour stage is a histogram aggregation.}
\end{table}

\phantomsection
\subsection*{Statistical tests and model selection\texorpdfstring{\appmark{app:nb}}{}}
\label{app:nb}
\paragraph{The problem: cross-validation folds are not independent.}
When comparing two models using $k$-fold cross-validation, we obtain one paired performance difference per fold, $d_1,\ldots,d_k$. A conventional paired $t$-test assumes these differences are independent and estimates the variance of their mean as $\hat{\sigma}^2/n$, where $n$ is the number of paired observations.

This assumption does not hold in cross-validation. Although the test folds are disjoint, their training sets overlap substantially: in our five-fold gene split of data0, 73.4\% of fold 0's training sites also appear in fold 1's training set. Models trained on such similar data are not independent, so their performance differences across folds may also be correlated. Treating these differences as independent underestimates uncertainty and can produce overly optimistic $p$-values.

\paragraph{The Nadeau--Bengio correction.}
The corrected resampled $t$-test \citep{nadeau2003} accounts approximately for this dependence by adding a term based on the ratio of test-set to training-set size. For $k$ equally sized folds, this ratio is $1/(k-1)$, giving
\[
\widehat{\operatorname{Var}}(\bar d)
=
\hat{\sigma}^2
\underbrace{\left(
\frac{1}{n}+\frac{1}{k-1}
\right)}_{\text{correction for overlapping training sets}},
\]
instead of the conventional $\hat{\sigma}^2/n$.

Here, $\hat{\sigma}^2$ is the sample variance of the paired differences, and $n$ is the total number of paired observations ($k$ for one cross-validation run, or $rk$ for $r$ repetitions). The additional term increases the estimated uncertainty and produces a more conservative significance test.

\paragraph{Example: comparing two feature representations.}
We compared two LightGBM feature sets on data0: \texttt{pooled\_v1}, which summarises each read feature by its per-site mean and standard deviation, and \texttt{quantiles\_v1}, which adds five per-site quantiles (5th to 95th percentile), the interquartile range and the tail spread of each feature. Both include the read count and the 7-mer, and both were evaluated on identical held-out gene folds. Every $p$-value below is one-sided, testing $H_0$: the two feature sets have equal PR AUC, against $H_1$: \texttt{quantiles\_v1} has the higher PR AUC. (Our harness reports two-sided values; because the observed difference favours quantiles, the one-sided value is exactly half.)

\begin{center}
\small
\begin{tabular}{lcc}
    \toprule
    \textbf{Statistical test}
    & \textbf{5 folds}
    & \textbf{10 repetitions $\times$ 5 folds} \\
    \midrule
    Unpaired Welch $t$-test
    & $p=0.15$ & -- \\
    Naive paired $t$-test
    & $p=0.023$ & $p=5.5\times10^{-21}$ \\
    Nadeau--Bengio corrected $t$-test
    & $p=0.065$ & $p=4.8\times10^{-5}$ \\
    \midrule
    Folds favouring quantiles
    & 5/5 & 50/50 \\
    PR AUC improvement
    & $+0.0148$ & $+0.0164$ \\
    \bottomrule
\end{tabular}
\end{center}

\paragraph{Interpretation.}
With five folds, the quantile representation outperforms the pooled representation on every fold. The naive paired test gives $p=0.023$, suggesting that quantiles are better at $\alpha=0.05$. However, after accounting for overlapping training sets, the corrected test gives $p=0.065$. Five consistent but correlated results are not equivalent to five independent experiments.

To obtain stronger evidence, we repeated five-fold cross-validation using ten independently seeded gene partitions, yielding 50 paired differences. Quantiles outperform pooled features in all 50 comparisons, with a mean PR AUC gain of 0.0164. The corrected test now gives $p=4.8\times10^{-5}$, providing strong evidence that the improvement is systematic rather than specific to one gene partition.

Importantly, repeating cross-validation does not eliminate dependence between training sets. For $k=5$, the corrected variance factor is $1/n+1/4$: increasing $n$ reduces the first term, but the overlap correction remains. Thus, repeated cross-validation provides more evidence across different partitions while the corrected test continues to account approximately for their dependence. The naive test, which ignores the overlap, becomes ever more overconfident as $n$ grows ($p=5.5\times10^{-21}$ at $n=50$).
% SOURCE: src/m6a/compare.py (docstring, corrected_variance); docs/decisions/0005
% (73.4% overlap, data0 split) and 0006 (W&B run y2lk7ilf: configs/quantiles.yaml
% --repeats 10 --compare-features pooled_v1; +0.0148 at 5 folds, naive p 1.1e-20
% at 50). Two-sided harness values (0.3095, 0.0465, 0.1305, 1.1e-20, 9.5e-5)
% halved to one-sided; src/m6a/features/{pooled,quantiles}.py (feature definitions).

\paragraph{Seeds, ties and model selection.}
Each design is trained with three seeds, and their within-file ranks are
averaged; differences below about 0.005, the observed seed spread, are treated
as ties. Models are ranked by their worse PR AUC gain over the baseline across
the two cell lines, with ties broken on data1's new genes
\citep{gulrajani2021,sagawa2020}. The rule guards against gains that hold in
one cell line only: an early feature set improved PR AUC by 0.065 within data0
but by only 0.016 on data1, and by $-0.003$ on genes unique to data1.
% SOURCE: docs/decisions/0029, 0032; findings.md section 1.

\phantomsection
\subsection*{Supplementary Methods: evaluation protocol\texorpdfstring{\appmark{app:evaluation-details}}{}}
\label{app:evaluation-details}


\paragraph{Datasets.}
We use two labelled datasets: data0 (HCT116; 121,838 candidate sites, 4.49\% modified) and data1 (HEK293T; 90,810 sites, 7.24\% modified), hereafter cell lines 1 and 2. data0 corresponds to SG-NEx HCT116 replicate 3, run 1, matched site by site. data2, synthetic RNA sequenced at known modification fractions, is reserved for external evaluation. Generalisation across cell lines is essential: an early feature set that performed well within data0 transferred poorly to data1, suggesting overfitting to HCT116-specific labels.

\paragraph{Gene-held-out and cross-cell-line evaluation.}
We use five-fold cross-validation grouped by gene, with a single gene partition shared by both cell lines. No gene, transcript, site or read from a test gene appears in training, preventing leakage through overlapping transcripts or neighbouring sites \citep{bbag434}. Within each training fold, 10\% of genes are reserved for early stopping, and label-derived parameters are fitted only on training genes.

We train each design under two settings: (i) cell line 1 only, to assess transfer to the unseen cell line 2; and (ii) both cell lines, matching final deployment. Both settings are evaluated on held-out genes from each cell line and on data1 genes absent from data0 (``new genes''). Results are out-of-fold unless explicitly marked in-sample. The submitted model is refitted on all labelled data.

\paragraph{Read-depth robustness and external validation.}
Although labelled sites have at least 20 reads, SG-NEx sites have a median of only 3--5. To assess robustness, we randomly subsample held-out sites and their neighbours to 1, 3 and 10 reads, evaluating models trained at full depth. We additionally evaluate on two independent SG-NEx HCT116 sequencing runs using data0's labels, and on the held-out synthetic data2. Neither external dataset is used for training.

\paragraph{Model comparison and selection.}
All comparisons are paired on identical held-out sites and reported as differences in PR AUC, with 95\% confidence intervals and one-sided $p$-values (testing whether the new design is better) from 1,000 resamples of whole genes. Where only fold-level scores are available, we use the Nadeau--Bengio corrected resampled $t$-test \citep{nadeau2003}. Because transfer to new cell lines is the priority, models are ranked by the smaller of their PR AUC gains over the baseline on the two cell lines.\applink{app:nb}
% SOURCE: docs/decisions/0029, 0032; analysis/representation/day_significance.py
% (gene bootstrap; p = fraction of resamples with the wrong sign, i.e. 1 - win
% rate); scripts/evaluate.py (Nadeau-Bengio); findings.md section 1;
% analysis/newdata/sgnex_source_match.py; cell lines from course staff (14 Oct).

\paragraph{Expanded model specification.}
\paragraph{Components.}
\begin{itemize}
    \item \textbf{Read encoder} (Section~\ref{sec:reads}): A DeepSet encoder
    applies a shared network $\phi$ to each read, pools the resulting
    embeddings using their mean, spread and log read count, and passes these
    summaries through a site network $\rho$. An auxiliary head scores each
    site from its own reads alone during training.

    \item \textbf{Two-head H2GCN} (Section~\ref{sec:graph}): Each candidate
    site is a graph node connected to sites on the same transcript within
    200\,nt. Two graph layers capture one- and two-hop neighbours. Four
    representations --- the site's own embedding, its one-hop and two-hop
    neighbourhoods, and a whole-transcript summary --- are kept separate
    before concatenation \citep{xu2018}.

    Two output heads learn the labels of the two cell lines independently,
    while sharing the underlying network. This accommodates disagreements
    between their labels; the two predictions are averaged at
    inference. This component uses a 200\,nt graph and a whole-transcript channel;
    the constrained variants use local distance-weighted context
    (Supplementary Methods).

    \item \textbf{Constrained variants} (Section~\ref{sec:constrain}):
    Residual + dropout and scalar messages + dropout. In the former,
    distance-weighted neighbour evidence adjusts the own-reads score; in the
    latter, neighbours transmit their own-reads score and read count rather
    than a full embedding. Both randomly hide neighbours during training.
\end{itemize}

\paragraph{Why two heads beat one pooled head.}
\label{app:twohead}
Both extra heads --- the own-reads head and the second cell line's head ---
are multi-task learning \citep{caruana1997}: one shared network, several
related targets.
Let $\operatorname{logit}P(Y^{(c)}=1\mid x)=w_c^\top h(x)+b_c$. If
$w_1=w_2$, both the pooled ideal
$\tfrac12(\sigma(w^\top h+b_1)+\sigma(w^\top h+b_2))$ and the head average
$w^\top h+\tfrac{b_1+b_2}{2}$ are monotone in $w^\top h$: same ranking. If
$w_1\neq w_2$ (our case), the pooled model fits one function to two
conflicting targets, while two heads resolve the conflict at the output.
Averaging logits is also not averaging probabilities: heads of 0.99 and 0.50
give 0.745 by probability and $\sigma(2.3)=0.909$ by logit.

\begin{table}[H]
    \centering\small
    \begin{tabular}{lccc}
        \toprule
        & \textbf{Cell line 2} & \textbf{Cell line 1} & \textbf{Seed SD (\cltwo)} \\
        \midrule
        One head (\mdl{h2gcn}) & +0.031 (+0.036, +0.027, +0.030) & +0.093 & 0.005 \\
        Two heads (\mdl{h2gcn_twohead}) & +0.039 (+0.039, +0.037, +0.041) & +0.108 & 0.002 \\
        \bottomrule
    \end{tabular}
    \caption{Same architecture, one vs two heads, trained on both cell lines,
    3 seeds; PR AUC gain over LightGBM.}
\end{table}

\phantomsection
\subsection*{Complementary components and intermediate-model overview\texorpdfstring{\appmark{app:model-overview}}{}}
\label{app:model-overview}
\label{sec:optim}
\paragraph{Combining Complementary Networks}
\paragraph{Why ensemble three networks?} The constrained variants transfer
best, but on cell line 1 they trail our 7 October intermediate submission (an
ensemble of the two-head network and H2GCN) by 0.019 and 0.038 PR AUC,
respectively. We therefore retain the two-head network's cell line 1
strength, since the test set may come from another HCT116 run, while
incorporating the two designs that generalise better.

Compared with the intermediate submission, the final ensemble performs
similarly on cell line 1 ($+0.001$, $p=0.39$), and improves on cell line 2
($+0.006$, $p<0.001$) and data1's new genes ($+0.012$,
$p=0.02$).~(Table~\ref{tab:ensembles}) The gains are small, and cell line 1 is a
tie.
% SOURCE: results/day_significance.csv (vs = "shipped ensemble"); decision 0034.

% Supporting findings remain deliberately expansive in this structural draft.
\paragraph{Overview of Intermediate Models}


Figure~\ref{fig:models_dist} places every model of Section~\ref{sec:dev} on
one scale, all trained on both cell lines and compared with
quantiles\,+\,LightGBM fold by fold. It retraces the development; the raw
values and the evidence for each step are collected in Table~\ref{tab:models18}.
\begin{itemize}
    \item Learned read encoders beat hand-made read summaries on both cell lines.
    \item A naive graph (GCN) is worse than no graph; separate channels for
    the site and its neighbours (H2GCN) make the graph help.
    \item More freedom (wider context, attention, two heads) helps cell line 1
    but not cell line 2.
    \item The constrained designs are the best single networks on cell line 2,
    and lose that edge when their neighbours are replaced by random ones.
    \item The final ensemble ranks first on both cell lines; no single network
    does.
\end{itemize}
% SOURCE: results/report_models_folds.csv (version=both), medians of
% "gain over lightgbm" over 15 fold x seed points.

\begin{figure}[H]
    \centering
    \reportfig[width=0.85\linewidth]{models_violin_std_both.png}
    \caption{PR AUC gain over quantiles\,+\,LightGBM on the same fold (trained on both cell lines), rescaled within each cell line so the two share a scale; 0 is LightGBM. Points are fold--seed pairs, bars medians, hollow markers controls.~(Table~\ref{tab:models18})}
    \label{fig:models_dist}
\end{figure}
% SOURCE: analysis/representation/report_models.py (paired_violins, standardise=True);
% unstandardised version: models_violin_both.png.
% results/report_models_folds.csv. Also available: models_violin_cl1.png
% (cell line 1 only arm), models_distribution_{cl1,both}.png (two-panel
% version), models_scatter_{pr,roc}_{cl1,both}.png.


\paragraph{Ensemble candidates.}
\begin{table}[H]
    \centering\small
    \begin{tabular}{lccc}
        \toprule
        \textbf{Candidate} & \textbf{Cell line 1} & \textbf{Cell line 2} & \textbf{data1 new genes} \\
        \midrule
        Residual + dropout alone & $-0.019$ & $-0.001$ & -- \\
        Scalar messages + dropout alone & $-0.038$ & +0.001 & -- \\
        Residual + scalar messages & $-0.019$ & +0.006 & -- \\
        Two-head + residual & +0.004 & $-0.000$ & -- \\
        Two-head + residual + scalar messages & $-0.001$ & +0.006 & +0.012 \\
        \textbf{Two-head + residual + scalar msg.\ + dropout} & \textbf{+0.001} & \textbf{+0.006} & \textbf{+0.012} \\
        \bottomrule
    \end{tabular}
    \caption{Every ensemble considered: paired PR AUC difference from the intermediate leaderboard
    ensemble (two-head + H2GCN), both-cell-lines arm, every component
    averaged over 3 seeds. Bold: the submitted model ([$-0.005$, +0.006],
    [+0.002, +0.009], [+0.000, +0.023]).}
    \label{tab:ensembles}
\end{table}
% SOURCE: results/day_significance.csv (vs = "shipped ensemble"); decision 0034.

\paragraph{The models compared.}
The evidence behind each takeaway of Section~\ref{sec:optim}, as median PR AUC
gains over quantiles\,+\,LightGBM (cell line 1 / cell line 2, trained on both;
Table~\ref{tab:models18}).
\begin{itemize}
    \item \textbf{Reads.} With no neighbours, DeepSet gains +0.056 / +0.028.
    \item \textbf{Naive vs separate-channel graph.} The GCN falls below
    LightGBM ($-0.133$ / $-0.099$); H2GCN gains +0.084 / +0.044.
    \item \textbf{Freedom.} A 400\,nt context, the graph transformer and the
    two-head network raise the cell line 1 gain to +0.10--0.13, but cell line
    2 stays at +0.04--0.05, no better than H2GCN.
    \item \textbf{Constraint.} The three constrained designs reach
    +0.057--0.061 on cell line 2, the best single networks there. With random
    neighbours in place of real ones, scalar messages drop from +0.061 to
    +0.032.
    \item \textbf{Ensemble.} The final ensemble reaches +0.134 / +0.068, ahead
    of the intermediate submission (+0.127 / +0.060). The best single network
    on cell line 1 (two-head, +0.128) and on cell line 2 (scalar messages,
    +0.061) are different networks.
\end{itemize}
\begin{table}[H]
    \centering\small
    \begin{tabularx}{\linewidth}{lccY}
        \toprule
        & \multicolumn{2}{c}{\textbf{Gain over LightGBM}} & \\
        \cmidrule(lr){2-3}
        \textbf{Model} & \textbf{\clone} & \textbf{\cltwo} & \textbf{Mechanism and lesson} \\
        \midrule
        Logistic regression & $-0.061$ & $-0.064$ & Per-site mean/sd: the required simple baseline \\
        Quantiles + LightGBM & 0 & 0 & Hand-made read-distribution features; the reference for every gain \\
        m6Anet (pretrained) & +0.032 & $-0.040$ & Multiple-instance learner; home advantage on \clone, fair test on \cltwo \\
        DeepSet & +0.056 & +0.028 & Learned read encoder, mean + spread pooling: learning from raw reads beats hand-made summaries \\
        Set Transformer & +0.056 & +0.021 & Attention between reads: no gain; reads are interchangeable molecules \\
        GCN & $-0.133$ & $-0.099$ & Site averaged with neighbours, which assumes homophily: dilutes the rare modified site, fails \\
        GAT & +0.068 & +0.008 & Attention over neighbours: helps \clone, barely \cltwo \\
        H2GCN + own-reads head & +0.084 & +0.044 & Separate self/neighbour channels: corroboration works \\
        H2GCN, 400\,nt & +0.112 & +0.038 & Wider context helps only the line it learns most from \\
        Graph transformer & +0.103 & +0.043 & Better within each cell line, not across: capacity learns that line's labelling \\
        Two-head H2GCN & +0.128 & +0.047 & One head per labelling; strongest single network on \clone \\
        Scalar messages & +0.082 & +0.061 & Barest corroboration; best single network on \cltwo \\
        \quad control: random neighbours & +0.057 & +0.032 & Real neighbours replaced by random ones: the gain collapses, so corroboration is the signal \\
        Residual + dropout & +0.114 & +0.057 & Constrained and robust; best all-rounder \\
        \quad control: no dropout & +0.107 & +0.051 & Without dropout both gains shrink: dropout is the active ingredient \\
        Scalar messages + dropout & +0.089 & +0.059 & Dropout lifts the scalar design on \clone \\
        Intermediate ensemble & +0.127 & +0.060 & Two-head + \mdl{h2gcn_aux}, submitted 7~Oct \\
        Final ensemble & +0.134 & +0.068 & Two-head + \mdl{res_gate} + \mdl{scalar_drop}: highest on both cell lines \\
        \bottomrule
    \end{tabularx}
    \caption{The 18 models in Figure~\ref{fig:models_dist}, with the raw
    values behind it: median PR AUC gain over quantiles\,+\,LightGBM on the
    same fold (15 fold $\times$ seed points; 5 for the logistic baseline and m6Anet), trained on both cell lines. Figure~\ref{fig:models_dist} divides each cell line's gains by the standard deviation of the model medians (0.070 on cell line 1, 0.047 on cell line 2).}
    \label{tab:models18}
\end{table}
% SOURCE: results/report_models_folds.csv (version=both), medians of
% "gain over lightgbm". The previous version of this table quoted gains from
% the cell-line-1-only arm (unseen cell line / cell line 1); those are in
% findings.md and report_models_pooled.csv (version=cl1).

\phantomsection
\subsection*{The constrained designs in detail\texorpdfstring{\appmark{app:designs}}{}}
\label{app:designs}
All three designs keep the read encoder and the own-reads auxiliary head of
Section~\ref{sec:graph}, and each was compared with H2GCN trained on the same
seeds. The batch's predicted outcomes were written down before it ran.

\paragraph{Distance weighting.} Neighbours out to 150\,nt are averaged with
weights $e^{-d/\lambda}$ (row-normalised), where $\lambda$ is fitted in each
fold to the decay of co-modification in cell line 1's training labels
(101.5--122.1\,nt), following continuous distance filters
\citep{schutt2017}. A variant averages neighbours separately in three bands
(0--30, 30--75, 75--150\,nt) and mixes the bands with fixed weights taken from
the training labels \citep{corso2020}; its control reverses those weights.

\paragraph{Residual + dropout.} The score is
$z_i=a_i+g_i\,\Delta_i\,\mathbbm{1}[\deg_i>0]$: the site's own-reads score
$a_i$ plus a neighbour correction $\Delta_i$ (from the distance-weighted
neighbours above), scaled by a gate
$g_i=\sigma(\alpha\log n_i+\beta\,|a_i|+c)$ that can give neighbours more
weight when the site has few reads $n_i$ or an uncertain own score. With no
neighbours, $z_i=a_i$ exactly. This follows Correct-and-Smooth, in which a
base prediction is corrected by its graph neighbours \citep{huang2021}. In
training, 30\% of neighbour links and 15\% of neighbour sites are dropped at
random \citep{rong2020,feng2020}. The gate learned little: trained on cell
line 1 alone it stayed at a constant $\approx$0.5; trained on both lines it
learned the intended direction, but weakly (from 0.36 at 20--30 reads to 0.31
at 100+ reads), because every training site has at least 20 reads and the
1--3-read regime where it would matter never appears in training.

\paragraph{Scalar messages.} Each neighbour passes only its own-reads score
and its read count. A small network combines the site's own score with three
summaries: the distance-weighted mean of its neighbours' scores, the highest
neighbour score within 75\,nt, and the neighbours' mean log read count. This
is the factorised message of the corroboration derivation (Eq.~\eqref{eq:factorised})
with the evidence term $\varphi$ learned. Its control replaces the
neighbours' scores with those of unrelated sites.
% SOURCE: analysis/representation/graph.py (forward_v2, comod_decay_scale,
% comod_band_weights), xsrc_nets.py (GRAPH_VARIANTS); results/probe_gate_params.csv;
% findings.md sections 12A, 12E.

\paragraph{Constrained design definitions.}
\paragraph{Three constrained designs.}
Each design builds one property of corroboration into the model rather than leaving it to be learned, and each was tested against a control that removes that one property.\applink{app:designs}
\begin{itemize}
    \item \textbf{Distance weighting.} Nearer neighbours count more, with the weight fading over about 100\,nt, as co-modification does (Section~\ref{sec:cluster}).
    \item \textbf{Residual + dropout.} The site's own-reads score is the starting point, and its neighbours may only adjust it; a site with no neighbours keeps its own score. During training, randomly chosen neighbours are hidden, so the model cannot depend on any one of them.
    \item \textbf{Scalar messages.} Each neighbour passes just one number --- how modified it looks from its own reads --- together with its read count, instead of a learned vector.
\end{itemize}

\phantomsection
\subsection*{Supplementary Methods: training and implementation\texorpdfstring{\appmark{app:training-details}}{}}
\label{app:training-details}
\paragraph{Training.} All networks use weighted binary cross-entropy, with
positive examples weighted by the negative-to-positive ratio ($\approx 21$).
Optimisation uses AdamW (learning rate $10^{-3}$, weight decay $10^{-4}$),
with warm-up, learning-rate halving after 8 epochs without validation
improvement, and early stopping after 25 epochs without improvement in
validation PR AUC.

Batches contain whole transcripts, preserving each site's full graph
neighbourhood. All evaluated networks stopped on early-stopping patience
rather than the time limit, so none was cut short. The six final networks
(two seeds per architecture) were fitted on every labelled site of both
files, apart from the 10\% of genes held out for early stopping.
% SOURCE: analysis/representation/graph.py (train), xsrc_nets.py (GRAPH_VARIANTS),
% convergence.py; docs/decisions/0033 (final_fit.py, same 10% validation genes), 0034.

\paragraph{Inference and output.} The six networks are exported to NumPy,
allowing \texttt{scripts/predict.py} to run without a deep-learning
framework. Because the networks produce scores on different scales, each
network ranks sites within the input file, and these ranks are averaged to
produce a final score in $[0,1]$.

The scores represent relative rankings, not calibrated probabilities. In
particular, positive-class weighting shifts logits by
$\log w \approx 3.06$. This does not affect the
ranking-based evaluation, but limits probabilistic interpretation
(Section~\ref{sec:limits}).

\paragraph{The class weight inflates scores by an exact amount.}
\label{app:classweight}
With positives up-weighted by $w=(1-\pi)/\pi\approx21.25$, the weighted loss
for a site with true probability $q$ is
$\mathcal{L}(p)=-wq\log p-(1-q)\log(1-p)$, minimised at
$p_w=\frac{wq}{wq+1-q}$, i.e.
\[
\operatorname{logit}p_w=\operatorname{logit}q+\log w,\qquad
\log w\approx3.06 .
\]
Ranking metrics are unaffected (a constant shift does not change the order),
but scores are not probabilities. Moving to a target with base rate $\pi_t$,
$\operatorname{logit}q_t=\operatorname{logit}q_s-\operatorname{logit}\pi_s+\operatorname{logit}\pi_t$
\citep{saerens2002}.

\phantomsection
\subsection*{Supplementary methods: comparators, read-representation screening and co-modification normalisation\texorpdfstring{\appmark{app:methods-extra}}{}}
\label{app:methods-extra}
\label{app:comparator-details}
\paragraph{Comparators.} We compare the final model with three others, all
scored on the same held-out sites: a simple baseline (logistic regression on
each site's per-feature mean and standard deviation plus its 7-mer);
quantiles\,+\,LightGBM \citep{ke2017} (per-site quantiles of each read
feature; our first submitted model and the reference for every gain in
Section~\ref{sec:dev}); and pretrained m6Anet \citep{hendra2022}, run through
its own pipeline. m6Anet was trained on SG-NEx HCT116 direct RNA with HCT116
m6ACE-seq labels \citep{koh2019} --- data0's cell line, assay and pipeline, and almost
certainly many of the same sites --- so its cell line 1 score is not held out.
We also retrained m6Anet on our gene folds of data0; it scored below the
pretrained model on the same sites (0.488 vs 0.503 PR AUC, 0.924 vs 0.932 ROC
AUC), so the pretrained model is the comparator.
% SOURCE: W&B runs m6anet_retrained_ourfolds (8iwst4yg, 25 Sep; oof over 5
% gene folds of data0, the earlier data0-only split, seed 4262) and the
% pretrained model on the same sites; analysis/m6anet/run_ronin.sh retrain/score.
% SOURCE: configs/baseline.yaml (pooled_v1 + logistic); configs/quantiles.yaml;
% analysis/m6anet/ (pretrained HCT116_RNA002); findings.md section 13.

\label{app:screen-protocol}
\paragraph{Representing a variable set of reads (read-representation screening).}
Each candidate site contains 20--991 reads in training, but often only 1--3 in SG-NEx. Each read represents a separate molecule with nine current features, while modification labels are available only at the site level. Since reads have no meaningful order, the site encoder must handle variable read counts and be \emph{permutation-invariant}. The Deep Sets framework \citep{zaheer2017} provides a natural approach: encode reads independently, pool their representations, and predict from the pooled vector.

Existing approaches include fixed per-site statistics, as in our LightGBM baseline, and learned per-read representations, as in m6Anet \citep{hendra2022} and DeepMod \citep{liu2019deepmod}. We therefore investigated whether learning a representation of the entire read set improves on independently encoding and summarising reads.

\paragraph{Representation screening.}
We screened 14 representations, including hand-crafted statistics, unsupervised encoders (autoencoders, an LSTM across the three ordered pore positions, a masked Set Transformer and contrastive learning), and label-supervised encoders (an autoencoder with a prediction head, attention MIL and DeepSet).

Each representation was scored by simple probes trained on its output, independently of the encoder's own prediction head (Table~\ref{tab:reads}).

\label{app:cluster-protocol}
\paragraph{Co-modification normalisation: measuring local clustering.}
Let $Y_i\in\{0,1\}$ indicate the modification status of site $i$, and $d_{ij}$ the nucleotide distance between two candidate sites on the same transcript. We first measure the co-modification ratio
\[
R(d)=
\frac{P(Y_i=1,Y_j=1\mid d_{ij}=d)}
{P(Y_i=1)P(Y_j=1)},
\]
which compares observed co-modification with that expected under independence. Within 25\,nt, $R(d)$ is 7.4 in cell line 1 but 4.5 in cell line 2. However, this measure also captures differences in modification prevalence \emph{between transcripts}: two sites on an m6A-rich transcript are more likely to be modified regardless of their separation.

To isolate \emph{local} clustering, we instead normalise by the expected co-modification when modified sites are randomly distributed among candidate positions \emph{within each transcript}, preserving its modification rate:
\[
r(d)=
\frac{P(Y_i=1,Y_j=1\mid d_{ij}=d)}
{P_{\mathrm{own}}(Y_i=1,Y_j=1)}.
\]
Under this transcript-specific null model, $r(d)=1$ indicates no local clustering, while $r(d)>1$ indicates enrichment of nearby modified pairs.

\phantomsection
\subsection*{SG-NEx and GLORI: evaluation protocol\texorpdfstring{\appmark{app:glori-protocol}}{}}
\label{app:glori-protocol}
Reads were pooled across each cell line's samples before scoring. The analysis
uses \texttt{score\_raw}, the average raw output of the six shipped networks,
not the submission's average within-file rank. Candidate sites were mapped
from transcript coordinates using Ensembl 91 (GRCh38) exon annotation and
matched by chromosome, position and strand. Isoform duplicates were collapsed
by retaining the transcript copy with the most reads. Sequence checks matched
97.2\% of candidate 7-mers to the reference cDNA. Exact-coordinate matches
covered 50.1\%, 53.3\%, 71.0\% and 74.4\% of the GLORI lists in A549, H9,
HepG2 and HCT116; shifting by one base gave at most 0.23\% matches. These
checks support the coordinate convention without proving that every site maps.

GLORI positives are the authors' published calls, without an additional
threshold. In A549 and H9 we used
\href{https://www.ncbi.nlm.nih.gov/geo/query/acc.cgi?acc=GSE290408}{GSE290408},
files \path{A549_GLORI3_hg38_SNuc_KDremoved} and
\path{H9_GLORI3_hg38_SNuc_KDremoved}; these lists filter sites persisting
after METTL3 knockdown. HepG2 and HCT116 use
\href{https://www.ncbi.nlm.nih.gov/geo/query/acc.cgi?acc=GSE301419}{GSE301419},
\path{GSM9083038_HepG2_sgCTRL.totalm6A.FDR} and
\path{GSM9083040_HCT116_sgCTRL.totalm6A.FDR}. The latter are control-guide
samples from cells expressing dCas13b-FTO; background effects of this
engineered system cannot be excluded. They have smaller call lists
(7,356 and 6,680 sites) than A549 and H9 (47,259 and 34,906).

Evaluation includes only DRACH candidates in genes with at least one matched
GLORI positive. Unlisted candidates are inferred negatives, not experimentally
verified unmodified sites. HCT116 training sites were removed. The resulting
sets contain 518,544 A549, 570,700 H9, 162,764 HepG2 and 99,576 HCT116 genomic
candidates. Missing GLORI calls and differences in cultures can affect both
absolute metrics and relative model comparisons; neither PR AUC nor precision
is asserted to be a universal lower bound on biological accuracy.

\phantomsection
\subsection*{Depth-adjusted gene-group analysis\texorpdfstring{\appmark{app:sgnex-gene-methods}}{}}
\label{app:sgnex-gene-methods}
For each of the seven lines, genomic candidates were split into eleven pooled
depth bands: 1, 2, 3--4, 5--9, 10--19, 20--31, 32--46, 47--83, 84--303,
304--599 and 600+ reads. Sites in each band's top 5\% by raw score were
marked high-scoring. Each gene's value is its high-scoring site proportion;
genes with fewer than three candidates are excluded. Class effects compare
equally weighted gene means against other protein-coding genes, with 95\%
intervals from 500 gene resamples. Long non-coding RNAs are compared with
protein-coding genes. This measure describes relative candidate enrichment,
not the absolute number of modified sites or modified molecules.

The independent GLORI analysis uses candidates with at least 10 SG-NEx reads
in GLORI-covered genes. A site's GLORI label minus its depth-band positive
rate is averaged per gene, then compared between gene groups. Model effects
and GLORI effects therefore have different definitions, despite both being
expressed in percentage points. A consistent finding requires the model
effect to have the same sign in all seven lines, with intervals excluding
zero in at least four; GLORI must agree in all four lines, with intervals
excluding zero in at least two. The 54 classes include transcription factors,
ribosomal proteins, mitochondrial genes, long non-coding RNAs and 50 MSigDB
hallmark sets (v2023.2). There is no multiplicity correction, and sequence
motif composition is not adjusted for.

\begin{table}[H]
\centering\footnotesize\setlength{\tabcolsep}{3pt}
\begin{tabular}{lrrl}
\toprule
Gene class & Model effect (pp) & GLORI effect (pp) & Consistency criterion \\
\midrule
TNF signalling via NF-kB & +2.9 to +4.4 & +1.7 to +4.9 & yes \\
transcription factors & +2.0 to +3.2 & +0.7 to +2.3 & yes \\
early oestrogen response & +1.1 to +2.3 & +0.3 to +4.0 & yes \\
Long non-coding RNAs & +0.5 to +1.4 & +3.3 to +5.6 & yes \\
mitotic spindle & -1.6 to -1.1 & -1.7 to -1.3 & yes \\
mTORC1 signalling & -3.0 to -1.9 & -1.1 to -0.5 & yes \\
E2F targets & -3.1 to -2.4 & -2.8 to -1.1 & yes \\
MYC targets (v1) & -3.9 to -3.1 & -2.5 to -1.7 & yes \\
protein secretion & -3.8 to -3.5 & -1.9 to -0.8 & yes \\
ribosomal proteins & -3.7 to -1.9 & -0.5 to +1.5 & no: GLORI disagrees \\
mitochondrially encoded & -5.7 to -5.4 & no GLORI calls & negative control \\
\bottomrule
\end{tabular}
\caption{Gene-group effects: ranges across seven model-scored lines and four GLORI lines. pp denotes percentage points; model and GLORI effects use different depth-adjusted measures. Ribosomal proteins do not meet the cross-assay criterion. Mitochondrial genes are evaluable in six model-scored lines and lack GLORI calls.}
\label{tab:sgnex-gene-effects}
\end{table}

The raw A549 ribosomal-protein score difference was +0.38, while the
depth-adjusted model effect ranged from $-3.7$ to $-1.9$ percentage points
across lines. These quantities are on different scales and should not be
compared numerically. Their change in direction illustrates confounding;
GLORI effects ranged from $-0.5$ to +1.5 points and do not establish depletion.
Mitochondrial genes lacked sufficient eligible data in HCT116 and had no
GLORI calls; their low model scores are descriptive, not independent
confirmation of absence of m6A.

\phantomsection
\subsection*{Read-representation screening and pooling findings\texorpdfstring{\appmark{app:read-findings}}{}}
\label{app:read-findings}
\paragraph{From Reads to a Site Representation}



We compared supervised and unsupervised read representations using the
screening protocol in Supplementary Methods.\applink{app:methods-extra}
Table~\ref{tab:reads} shows representative probe results.

\begin{table}[H]
    \centering\small
    \begin{tabularx}{\linewidth}{lYcc}
        \toprule
        \textbf{Representation} & \textbf{Read aggregation} & \textbf{LightGBM} & \textbf{Logistic} \\
        \midrule
        LSTM autoencoder & Encodes three pore positions per read; unsupervised reconstruction & 0.295 & 0.259 \\
        Autoencoder + prediction & Per-read four-dimensional code with site-level supervision & 0.462 & 0.449 \\
        \textbf{DeepSet} & \textbf{Supervised per-read encoder with mean and spread pooling} & \textbf{0.479} & \textbf{0.509} \\
        \bottomrule
    \end{tabularx}
    \caption{PR AUC of probes trained on three of the 14 screened representations (held-out genes, cell line 1). Probes see genes the encoder did not; scores are not comparable with end-to-end models.~(Table~\ref{tab:reads_full})}
    \label{tab:reads}
\end{table}
% SOURCE: analysis/representation/results/summary.md (E1); PLAN.md (protocol).
% The screen's hand-made reference (quantiles_all_v1, 0.535) is deliberately not
% shown: it includes neighbour and transcript-position columns, so it is not a
% per-site comparison. The like-for-like comparison is end to end (Finding 3).

\paragraph{Finding 1: supervision matters more than encoder architecture.}
Supervised encoders consistently outperformed unsupervised ones. Adding a site-level prediction head to an autoencoder increased PR AUC from 0.228 to 0.462, while the unsupervised LSTM scored only 0.295. This suggests that reconstruction objectives primarily capture sequence-dependent current variation rather than the weaker m6A signal.

Figure~\ref{fig:read_umap} supports this interpretation: hand-crafted and unsupervised representations cluster mainly by sequence motif, whereas supervised representations separate sites by modification status. DeepSet learns an approximately one-dimensional modification axis (effective dimensionality 1.8 of 32), visible even after two-dimensional projection.

\begin{figure}[H]
    \centering
    \reportfig[width=\linewidth]{fig_read_umap.png}
    \caption{UMAP projections of four site representations, coloured by modification label (top) and motif-specific modification rate (bottom). Titles report 30-nearest-neighbour PR AUC in the original representation and its two-dimensional projection. Every screened representation is mapped in Figure~\ref{fig:read_umap_full}.}
    \label{fig:read_umap}
\end{figure}
% SOURCE: analysis/representation/report_umap.py reads; summary.md E4 (effective dims).

DeepSet also slightly outperformed the supervised autoencoder, both in the screen and when its features were added to the hand-crafted ones. Unlike the autoencoder, which compresses individual reads before summarising their codes, DeepSet jointly learns the read encoder and set-level aggregation from site labels. It also avoids a reconstruction objective and can be trained end to end within the subsequent graph model, allowing a site's prediction loss to update its neighbours' read encoders.
% SOURCE: results/r2.md (50 observations, cross-fitted, Nadeau-Bengio corrected).

\paragraph{Finding 2: simple pooling is sufficient.}
We next tested whether explicitly weighting or modelling interactions between reads improves DeepSet. A predefined acceptance criterion required a mean gain of at least +0.005 PR AUC and a win on at least four of five folds. Attention MIL \citep{ilse2018} ($-0.003$, $p=0.55$), a Set Transformer \citep{lee2019} ($-0.006$, $p=0.68$), and a residual DeepSet with 24 times more parameters (+0.006, $p=0.38$) all failed this criterion. On the unseen cell line, DeepSet also slightly outperformed the Set Transformer (0.395 vs 0.390).

All pooling variants learned a similar low-dimensional modification signal. This suggests that the mean and spread of read-level currents capture most of the useful evidence. More complex pooling may identify individual modified reads, but site-level labels cannot directly supervise this. Similarly, m6Anet found no significant benefit from attention pooling \citep{hendra2022}.
% SOURCE: results/diagnose2.md and r3_gate.md (round 3, the pre-registered gate);
% results/report_models_pooled.csv (cl1 arm, 3 seeds).

\paragraph{Finding 3: learned representations generalise and tolerate low depth.}
End-to-end DeepSet outperformed quantiles + LightGBM on cell line 1 (0.531 vs 0.476 PR AUC) and the unseen cell line (0.395 vs 0.372), with a significant cross-cell-line gain in the paired analysis (+0.017, 95\% CI [+0.009, +0.025], $p<0.001$).

Its advantage increased as read depth decreased. On cell line 2, with each site subsampled to 1 or 3 reads, the learned read encoder (evaluated by removing neighbours from the graph model) achieved PR AUCs of 0.225 and 0.269, compared with 0.136 and 0.200 for LightGBM. Thus, learned read features remain useful even when a site's evidence is sparse.
% SOURCE: results/report_models_pooled.csv (cl1 arm, 3 seeds);
% results/deepset_nn__xsrc.json (seed-0 interval); depth: findings.md section 8
% (depth asymmetry, "no neighbours" column) and results/depth_ensembles.csv
% (LightGBM, cell line 2, trained on both).

\paragraph{Design implication.}
We adopt \textbf{DeepSet} as the read encoder for all subsequent models: a shared per-read network followed by mean and spread pooling and the log read count. This representation is permutation-invariant, supports variable read depth, generalises better to unseen cell lines, and integrates naturally into end-to-end graph learning. No recurrent layer is needed, since the ordered pore positions within each read are already represented explicitly by the input features and 7-mer.


\paragraph{Read-encoder screen, all representations.}
\begin{table}[H]
    \centering\footnotesize\setlength{\tabcolsep}{4pt}
    \begin{tabularx}{\linewidth}{Ylcccc}
        \toprule
        \textbf{Representation} & \textbf{Code name} & \textbf{Logistic} & \textbf{$k$NN} & \textbf{MLP} & \textbf{LightGBM} \\
        \midrule
        \multicolumn{6}{l}{\emph{Hand-made features and controls}} \\
        Full hand-engineered set (incl.\ neighbour and transcript-position columns) & \mdl{hand} & 0.496 & 0.372 & 0.506 & 0.535 \\
        The same, PCA to 32 dimensions & \mdl{hand_pca} & 0.252 & 0.325 & 0.433 & 0.419 \\
        Network on the hand-made features (labels) & \mdl{mlp_hand} & 0.468 & 0.454 & 0.465 & 0.452 \\
        Currents normalised by their 7-mer's mean, no network & \mdl{kmer_norm} & 0.289 & 0.325 & 0.356 & 0.373 \\
        Untrained per-read network, pooled (control) & \mdl{random} & 0.165 & 0.202 & 0.191 & 0.264 \\
        \midrule
        \multicolumn{6}{l}{\emph{Encoders trained without labels}} \\
        Per-read autoencoder, 4-number code & \mdl{read_ae} & 0.194 & 0.215 & 0.211 & 0.228 \\
        LSTM across a read's 3 pore positions, reconstructs the read & \mdl{pos_lstm_ae} & 0.259 & 0.287 & 0.301 & 0.295 \\
        Self-attention across the 3 pore positions, reconstructs the read & \mdl{pos_attn_ae} & 0.227 & 0.242 & 0.273 & 0.264 \\
        Masked set transformer: predicts hidden reads from the others & \mdl{set_masked} & 0.451 & 0.380 & 0.478 & 0.443 \\
        Contrastive: two random halves of a site's reads agree & \mdl{contrastive} & 0.355 & 0.327 & 0.440 & 0.419 \\
        \midrule
        \multicolumn{6}{l}{\emph{Encoders trained with labels}} \\
        Autoencoder + site-level prediction head & \mdl{read_ae_pred} & 0.449 & 0.469 & 0.471 & 0.462 \\
        Gated-attention MIL & \mdl{attn_mil} & 0.502 & 0.493 & 0.494 & 0.470 \\
        \textbf{DeepSet: per-read network, mean + spread pooling} & \mdl{deepset} & \textbf{0.509} & \textbf{0.498} & \textbf{0.476} & \textbf{0.479} \\
        DeepSet without the 7-mer & \mdl{deepset_nokmer} & 0.502 & 0.486 & 0.500 & 0.463 \\
        \bottomrule
    \end{tabularx}
    \caption{The full read-encoder screen behind Table~\ref{tab:reads}: PR AUC
    of four probes trained on each representation, held-out genes of cell
    line 1 (5 folds; encoders trained on half of each training fold, probes
    on the other half). The hand-made row includes columns computed from
    neighbouring sites and transcript position, so it is not a per-site
    comparison; the like-for-like comparison with hand-made per-site features
    is end to end (Section~\ref{sec:reads}, Finding 3).}
    \label{tab:reads_full}
\end{table}
% SOURCE: analysis/representation/results/summary.md (E1); PLAN.md (encoders, protocol).

\paragraph{UMAP of every screened representation.}
Figure~\ref{fig:read_umap} shows four of these maps; the full set follows,
one row per representation, in the order of the code names below. Each row
shows the same held-out sites (fold 0: every modified site and 5,000
unmodified ones) coloured by label (positives in red), by the motif's
positive rate, and by log read depth, with the 30-nearest-neighbour PR AUC on
the full vector and on the 2-D map. A wider DeepSet (\mdl{deepset_large}),
not in the screen table, is included.

\begin{figure}[H]
    \centering
    \reportfig[width=0.88\linewidth]{fig_read_umap_full_1.png}
    \caption{UMAP of every representation, part 1 of 3: \mdl{hand},
    \mdl{hand_pca}, \mdl{random}, \mdl{mlp_hand}, \mdl{deepset}.}
    \label{fig:read_umap_full}
\end{figure}

\begin{figure}[H]
    \centering
    \reportfig[width=0.88\linewidth]{fig_read_umap_full_2.png}
    \caption{UMAP of every representation, part 2 of 3: \mdl{deepset_large},
    \mdl{read_ae}, \mdl{read_ae_pred}, \mdl{pos_lstm_ae},
    \mdl{pos_attn_ae}.}
\end{figure}

\begin{figure}[H]
    \centering
    \reportfig[width=0.88\linewidth]{fig_read_umap_full_3.png}
    \caption{UMAP of every representation, part 3 of 3: \mdl{set_masked},
    \mdl{attn_mil}, \mdl{contrastive}, \mdl{deepset_nokmer},
    \mdl{kmer_norm}.}
\end{figure}
% SOURCE: analysis/representation/figures/umap.png (visualise.py), cut into three
% by analysis/representation/report_umap_full.py; kNN numbers in
% results/projection_stats.json.

\paragraph{Read pooling.}
\label{app:pooling}
Figure~\ref{fig:pooling_umap} repeats the view of Figure~\ref{fig:read_umap}
for four ways of pooling reads, all trained on the labels. Every variant
learns the same single ``how modified'' axis, and nearest-neighbour PR AUC
differs by at most 0.03. Across 50 cross-fitted comparisons, a wider DeepSet (+0.020), the supervised autoencoder (+0.020), attention MIL (+0.022) and DeepSet (+0.022) add about the same to the hand-made features (each corrected $p<0.0001$). In the screen itself, DeepSet scored 0.479 against the autoencoder's 0.462 (Table~\ref{tab:reads_full}). Three later variants inside
the full model also failed to beat mean and spread: quantiles of the read
encodings ($-0.009$), the hand-made features fed in alongside the reads
($-0.021$), and a graph linking each read to its four most similar reads,
which was no better than the same graph with its links shuffled.

\begin{figure}[H]
    \centering
    \reportfig[width=\linewidth]{fig_pooling_umap.png}
    \caption{Four ways of pooling reads, all trained on labels, in the view
    of Figure~\ref{fig:read_umap}. The choice of pooling barely matters, so
    we use the simplest.}
    \label{fig:pooling_umap}
\end{figure}
% SOURCE: analysis/representation/report_umap.py pooling (attn_mil and
% deepset_large fold-0 vectors regenerated on Ronin, 13 Oct); results/r2.md;
% findings.md section 3 (later variants).

\phantomsection
\subsection*{Local clustering: supporting findings\texorpdfstring{\appmark{app:cluster-findings}}{}}
\label{app:cluster-findings}
\paragraph{Modified Sites Cluster Along a Transcript}


\paragraph{Motivation.}
A read encoder predicts each site independently, potentially discarding information from neighbouring sites. Previous single-molecule studies have found that nearby m6A sites tend to be modified together, particularly within 200\,nt \citep{m6aiso2025,kang2025}. We test whether a similar spatial pattern exists in our site-level labels and whether it transfers across cell lines.

Before transcript-specific normalisation, the co-modification ratio within
25\,nt was 7.4 in cell line 1 and 4.5 in cell line 2. To distinguish local
clustering from transcript-wide enrichment, we normalised against random
redistribution of modified labels within each transcript (Supplementary
Methods).\applink{app:methods-extra}

\begin{figure}[H]
    \centering
    \reportfig[width=0.85\linewidth]{fig_comodification.png}
    \caption{Co-modification by distance between candidate sites on the same transcript: raw ratio $R(d)$ (dotted) and transcript-normalised ratio $r(d)$ (solid; 95\% intervals from resampling transcripts).}
    \label{fig:comod}
\end{figure}

\begin{center}\small
\begin{tabular}{lcccccc}
    \toprule
    \textbf{Distance (nt)} & 1--25 & 25--50 & 50--100 & 100--200 & 200--400 & 400--800 \\
    \midrule
    $r(d)$, cell line 1 & 2.20 & 2.10 & 1.88 & 1.46 & 1.11 & 0.80 \\
    $r(d)$, cell line 2 & 2.13 & 2.00 & 1.73 & 1.34 & 1.02 & 0.74 \\
    \bottomrule
\end{tabular}
\end{center}

\paragraph{Finding: local clustering is consistent across cell lines.}
After normalisation, both cell lines show nearly identical patterns (Figure~\ref{fig:comod}). Sites within 25\,nt are co-modified approximately 2.2 times more often than expected from their transcript's modification rate (2.20 and 2.13, with overlapping 95\% intervals). This enrichment decays approximately exponentially, with a characteristic scale of 100--120\,nt, approaching the null expectation by 200--400\,nt. At larger distances, $r(d)$ falls below 1, consistent with modified pairs being concentrated locally rather than spread uniformly along transcripts.

Thus, while cell lines differ in \emph{which transcripts} are enriched for m6A, they exhibit a similar \emph{local clustering pattern}. This suggests that distance-dependent relationships between neighbouring sites may generalise better than transcript-level modification patterns. Although this analysis used all labelled sites to motivate the architecture, subsequent design choices are evaluated on held-out genes, and the model's distance scale $\lambda$ is fitted exclusively on training labels within each fold.
% SOURCE: analysis/newdata/distance_bands/bands.csv (distance_bands.py):
% R(d) = `total`, r(d) = `local_part` with local_lo/local_hi; intervals
% 2.20 [2.04, 2.38] vs 2.13 [1.96, 2.29] below 25 nt; decay scale from
% analysis/newdata/decay_scale/decay_scale.csv. report/theory.md section 1.

\paragraph{Implication for the model.}
Neighbouring sites may provide corroborating evidence: a neighbour with strong modification evidence should increase a site's predicted score, while evidence against modification should decrease it, with influence diminishing over distance. A Bayesian derivation formalises this intuition. The following analyses examine how to incorporate this evidence (Section~\ref{sec:graph}), how far it should extend (Section~\ref{sec:reach}), and how strongly its influence should be constrained (Sections~\ref{sec:capacity}--\ref{sec:constrain}).


\paragraph{The corroboration message and four design rules.}
\label{app:message}
These are idealised models: they assume a site's reads depend only on its own
modification state and, where stated, that neighbours are independent. The
trained networks do not compute these formulas literally; read them as the
principle the architecture implements.

If modified sites cluster (Section~\ref{sec:cluster}), evidence measured at one
site is evidence about its neighbours. Assuming a site's reads depend only on
its own state, Bayes' rule splits the log-odds that site $i$ is modified,
given its reads $E_i$ and a neighbour's reads $E_j$, into three additive terms:
\begin{equation}
\operatorname{logit}P(Y_i{=}1\mid E_i,E_j)=\underbrace{\operatorname{logit}\pi}_{\text{prior}}
+\underbrace{\log L_i}_{\text{own reads}}+\underbrace{m_{ij}}_{\text{neighbour}},\qquad
m_{ij}=\log\frac{r(d_{ij})\,\pi\,L_j+1-r(d_{ij})\,\pi}{\pi' L_j+1-\pi'},
\label{eq:message}
\end{equation}
where $\pi$ is the base rate, $L_i$ and $L_j$ are the likelihood ratios from
each site's own reads, $r(d)$ is the normalised (local) co-modification ratio at distance
$d$ (Section~\ref{sec:cluster}), and $\pi'=\pi(1-r\pi)/(1-\pi)$. Linearising in
$r-1$, with $r(d)-1\approx Ae^{-d/\lambda}$, the message factorises into a
distance term and an evidence term:
\begin{equation}
m_{ij}\approx A\,e^{-d_{ij}/\lambda}\,\varphi(L_j),\qquad
\varphi(L)=\frac{\pi(L-1)}{(1-\pi)\bigl(1+\pi(L-1)\bigr)}.
\label{eq:factorised}
\end{equation}
Four properties follow, each a rule for the model:
\begin{enumerate}
    \item[(i)] \textbf{Keep own and neighbour evidence separate.} They enter
    as separate terms; a site's own evidence keeps its full weight however
    many neighbours it has.
    \item[(ii)] \textbf{Neighbours confirm or contradict.} $m_{ij}$ is
    positive when the neighbour looks modified ($L_j>1$) and negative when
    it does not; a neighbour with no evidence ($L_j=1$) sends nothing.
    \item[(iii)] \textbf{Only nearby neighbours count.} Where $r=1$ the
    message is zero, so neighbours help only within the
    $\approx$100--200\,nt where co-modification exists.
    \item[(iv)] \textbf{Neighbours adjust, never override.} The message is
    bounded by $\log r\le\log2.2\approx0.79$ logits per neighbour, and its
    share of the evidence grows as the site's own reads shrink.
\end{enumerate}

The derivation of Eq.~\eqref{eq:message} follows.

\paragraph{Derivation.} With $Y_i\in\{0,1\}$, base rate $\pi$,
own-reads likelihood ratio $L_i=P(E_i\mid Y_i{=}1)/P(E_i\mid Y_i{=}0)$, and
assumption A1 ($E_i\perp E_j\mid Y_i$; $E_j\perp Y_i\mid Y_j$), Bayes' rule
gives
$\operatorname{logit}P(Y_i{=}1\mid E_i,E_j)=\operatorname{logit}\pi+\log L_i+m_{ij}$
with $m_{ij}=\log\frac{P(E_j\mid Y_i=1)}{P(E_j\mid Y_i=0)}$. Expanding over the
neighbour's state, $P(E_j\mid Y_i{=}y)=\sum_{y_j}P(E_j\mid Y_j{=}y_j)P(Y_j{=}y_j\mid Y_i{=}y)$,
with $P(Y_j{=}1\mid Y_i{=}1)=r\pi$ and
$P(Y_j{=}1\mid Y_i{=}0)=\pi(1-r\pi)/(1-\pi)\equiv\pi'$, and dividing by
$P(E_j\mid Y_j{=}0)$ gives Eq.~\eqref{eq:message}. As $L_j\to\infty$,
$m_{ij}\to\log r(d)$ (at most $\log 2.2\approx0.79$ logits below 25\,nt); as
$L_j\to0$, $m_{ij}\to\log\frac{1-r\pi}{1-\pi'}\approx-(r-1)\pi$, small when
$\pi$ is small. Differentiating at $r=1$ gives Eq.~\eqref{eq:factorised}, with
$\varphi$ increasing, $\varphi(1)=0$ and
$\varphi\in\bigl(-\tfrac{\pi}{1-\pi},\tfrac{1}{1-\pi}\bigr)$.

\paragraph{Several neighbours.} If neighbours were independent given $Y_i$
their messages would add. They are not --- neighbours co-modify with each
other --- so a sum would treat correlated evidence as independent and be
overconfident on dense transcripts. Our designs use a distance-weighted
\emph{average} (plus a has-neighbour flag and the best neighbour), the
conservative end.

\paragraph{Why GCN fails.} A GCN layer
$h_i'=\phi\bigl(W\frac{h_i+\sum_j h_j}{\deg_i+1}\bigr)$ gives a site with 4
neighbours one fifth of the weight on its own evidence, whereas in the
posterior $\log L_i$ keeps its full weight however many neighbours there are.
Averaging is also a poor fit because the graph is not homophilous for the rare
class: within 25\,nt, a modified site's neighbour is itself modified with
probability about $R(d)\,\pi$, that is $7.4\times4.5\%$ and $4.5\times7.2\%$,
both $\approx0.33$ in the two cell lines.

\phantomsection
\subsection*{Corroboration controls and depth asymmetry\texorpdfstring{\appmark{app:graph-findings}}{}}
\label{app:graph-findings}
\paragraph{A Site Graph Improves Prediction Through Corroboration}


\paragraph{Graph construction and architecture.}
Each candidate site is a node connected to the sites on the same transcript within 50\,nt. Graph layers combine its DeepSet representation (Section~\ref{sec:reads}) with its neighbours' representations to predict modification.

Crucially, \emph{the site's own evidence must remain separate from its neighbours'}. Standard graph convolutional networks (GCNs; \citealp{kipf2017}) assume \emph{homophily} --- that linked nodes tend to share a label --- and average the representations of connected nodes. However, modified sites are rare: even within 25\,nt, only about one in three neighbours of a modified site is also modified. Averaging therefore dilutes modification evidence; with four neighbours, the site's own reads retain only one-fifth of the weight. Accordingly, a GCN performed worse than the graph-free LightGBM baseline, by 0.14 and 0.10 PR AUC on cell lines 1 and 2.

H2GCN \citep{zhu2020}, designed for graphs without homophily, instead maintains separate channels for the site's own representation and its neighbours' mean. This improved PR AUC by 0.010--0.016 over DeepSet alone, and an auxiliary head that predicts each site from its own reads improved it further (on cell line 2, from +0.033 to +0.044 over LightGBM).
% SOURCE: findings.md section 3 (GCN; h2gcn vs deepset; h2gcn_aux vs h2gcn_local).

\paragraph{Why the graph helps: corroboration.}
We interpret the graph's benefit as using neighbouring sites as evidence of modification, rather than simply combining more features. Four analyses support this interpretation. The gain comes mainly from distinguishing sites in modified neighbourhoods from the rest, not from better predictions within either group.~(Table~\ref{tab:status}) The model depends on the correct neighbours and their sequences: replacing them with random sites removes most of the gain.~(Table~\ref{tab:swap}) The gain rises with the number of neighbours and does not increase with transcript size, which argues against a coverage artefact.~(Figure~\ref{fig:nbrcount}) Finally, removing neighbours within 20\,nt has almost no effect, so the graph does not simply re-measure the target site.
% SOURCE: results/gain_decomposition.csv (neighbour status, neighbour count),
% day_summary.md (scalar_msg vs scalar_random, 3 seeds), ablate_graph.csv (swaps),
% band_ablation.csv (0-10 / 10-20 nt).

Together, these results support \emph{corroboration}: neighbours whose reads suggest modification raise the target site's score, while neighbours suggesting no modification lower it. In the final model, a neighbour scoring 90\% raises a target scoring 20\% to about 45\%, and one scoring 5\% lowers it to about 12\% (Figure~\ref{fig:corroboration}).

\begin{figure}[H]
    \centering
    \reportfig[width=0.7\linewidth]{fig_corroboration_curve.png}
    \caption{Corroboration learned by the final scalar-message model, computed from its weights: the target's score against its neighbour's own-reads score, for targets scoring 5\%, 20\% and 50\% alone (dotted: no neighbour). Both sites have 30 reads, within 75\,nt; scores are relative, not calibrated.}
    \label{fig:corroboration}
\end{figure}
% SOURCE: models/final/scalar_drop_s{0,1}.npz (report_figures.py corroboration_curve).

\paragraph{Neighbours help most at low read depth.}
When a site has fewer reads, its own evidence weakens, making neighbouring evidence more valuable. On cell line 2, thinning a target to one read while retaining full-depth neighbours yields 0.367 PR AUC, compared with 0.225 without neighbours, approaching the 0.405 achieved by a full-depth site without neighbours. When neighbours are also sparse, as in low-coverage SG-NEx data, the benefit is smaller but remains positive (+0.030 PR AUC at three reads).~(Table~\ref{tab:swap})
% SOURCE: results/isolated_rescore.csv (findings.md section 8, depth asymmetry).

\paragraph{Design implication.}
All subsequent graph models preserve separate representations for a site's own reads and its neighbours, with an auxiliary own-read prediction head. This allows the model to use neighbours for corroboration without sacrificing the site's direct evidence.


Neighbour replacement and removal test reliance on context in a trained network. Such interventions can also change the input distribution; they do not by themselves establish a causal biological mechanism or replace a separately trained graph-free comparator.


\paragraph{Graph gain by neighbour status.}
\begin{table}[H]
    \centering\small
    \begin{tabular}{lrcc}
        \toprule
        \textbf{Neighbours within 50\,nt} & \textbf{Sites} & \textbf{Positive rate} & \textbf{Gain over DeepSet} \\
        \midrule
        At least one modified & 9,308 & 26.4\% & $-0.007$ [$-0.018$, +0.007] \\
        Only unmodified & 64,238 & 4.7\% & +0.013 [+0.001, +0.027] \\
        None & 16,906 & 6.3\% & +0.003 [$-0.011$, +0.019] \\
        All sites & 90,452 & 7.3\% & +0.024 [+0.015, +0.032] \\
        \bottomrule
    \end{tabular}
    \caption{\mdl{h2gcn_aux} (trained on both, scored on cell line 2): the gain
    is $\approx$0 within each neighbour-status group and comes from telling
    the groups apart.}
    \label{tab:status}
\end{table}

\paragraph{What neighbours provide, and depth asymmetry.}
Removing neighbours costs 0.030 PR AUC, but substituting a random site's reads
and 7-mer costs 0.126, and pairing the neighbours' real reads with random
7-mers costs 0.163. This suggests the model reads each neighbour's current
relative to the signal expected for its sequence, rather than the current
values alone. A trained control that replaces neighbours' scores with those of
unrelated sites gains only +0.017 on the unseen cell line, against +0.055 with
the real ones (Table~\ref{tab:batch}).
\begin{table}[H]
    \centering\small
    \begin{tabular}{lcc}
        \toprule
        \textbf{Neighbours given (scored site unchanged)} & \textbf{PR AUC (\cltwo)} & \textbf{Change} \\
        \midrule
        Real inputs (as trained) & 0.436 & -- \\
        Nothing & 0.405 & $-0.030$ \\
        Reads with their transcript's average removed & 0.388 & $-0.048$ \\
        Reads borrowed from same-5-mer sites elsewhere & 0.323 & $-0.113$ \\
        A random site's reads and 7-mer & 0.310 & $-0.126$ \\
        Real reads, random 7-mer & 0.273 & $-0.163$ \\
        \bottomrule
    \end{tabular}
    \quad
    \begin{tabular}{lccc}
        \toprule
        \textbf{Own reads} & \textbf{Nbrs full} & \textbf{Nbrs thinned} & \textbf{No nbrs} \\
        \midrule
        1 & 0.367 & 0.243 & 0.225 \\
        3 & 0.381 & 0.299 & 0.269 \\
        10 & 0.408 & 0.372 & 0.338 \\
        all & 0.436 & 0.436 & 0.405 \\
        \bottomrule
    \end{tabular}
    \caption{Left: swapping what neighbours provide. Right: depth asymmetry
    --- the scored site thinned, with neighbours at full depth, equally
    thinned, or absent. \mdl{h2gcn_aux}, trained on both, cell line 2.}
    \label{tab:swap}
\end{table}
% SOURCE: results/ablate_graph.csv, isolated_rescore.csv

\paragraph{Graph gain by number of neighbours.}
\begin{figure}[H]
    \centering
    \reportfig[width=0.7\linewidth]{fig_neighbour_count.png}
    \caption{Graph gain over DeepSet (the same read encoder, no graph) by the
    number of other candidate sites within 50\,nt, with 95\% gene-resampled
    intervals; trained on cell line 1. The gain rises with neighbours on cell
    line 1 and is smaller and flatter on cell line 2. It does not rise
    steadily with transcript size, a proxy for coverage, which argues
    against a coverage artefact.}
    \label{fig:nbrcount}
\end{figure}
% SOURCE: results/gain_decomposition.csv.

\paragraph{Where the key assumption fails.}
\label{app:assumption}
The derivation treats a neighbour's reads as separate evidence, which holds
only if a site's reads depend on its own modification state alone. That fails
at very short range: one m6A distorts the current over about $\pm$10\,nt
\citep{kang2025}, so a ``neighbour'' within $\approx$20\,nt partly re-measures
the site itself, and a model would count one piece of evidence twice. The
model's gain does not come from there: removing every neighbour within 0--10
or 10--20\,nt changes PR AUC by at most 0.001 (Section~\ref{sec:graph}).

\phantomsection
\subsection*{Reach and capacity: supporting findings\texorpdfstring{\appmark{app:reach-capacity}}{}}
\label{app:reach-capacity}
\label{app:radius-findings}
\paragraph{How Far Should Neighbours Reach?}


\paragraph{About 100\,nt transfers; wider context does not.} Local
co-modification fades by about 200\,nt (Section~\ref{sec:cluster}), so
neighbours further away should add nothing that holds across cell lines. We
changed only the radius of the graph (Figure~\ref{fig:radius}; full evaluations in Table~\ref{tab:radius}).
On the training cell line, wider was always better: the gain over LightGBM
rose from +0.081 at 50\,nt to +0.112 at 400\,nt. On the unseen cell line it
peaked at 100\,nt (+0.043) and fell back by 200--400\,nt (+0.032--0.034).
Removing one distance band at a time from a trained 400\,nt network shows the
same split: nearby neighbours help both cell lines, while neighbours
200--400\,nt away help only the line the model was trained on ($-0.014$ vs
+0.001 when removed).~(Table~\ref{tab:bands})
% SOURCE: results/h2gcn_aux_r*_nn__xsrc.json (radius); results/band_ablation.csv.

\begin{figure}[H]
    \centering
    \reportfig[width=0.6\linewidth]{fig_radius.png}
    \caption{PR AUC gain over LightGBM as the neighbour radius widens, with only the radius changed. Blue: cell line 1, trained on both lines; orange: cell line 2, trained on cell line 1 only (unseen).}
    \label{fig:radius}
\end{figure}

\paragraph{Why the wider context does not transfer.} Within about 100\,nt, the
normalised ratio $r(d)$ shows clustering that both cell lines share. Further
out, a neighbour average mostly reflects how m6A-rich the whole transcript
is --- the transcript-level enrichment that separates the two raw ratios
$R(d)$ in Figure~\ref{fig:comod} --- and that differs between the cell lines.
The shared part follows gene structure: in both cell lines, m6A is depleted near exon junctions and enriched further into exons and after the stop codon.~(Figure~\ref{fig:genestructure}) This agrees with the exon-junction-complex model, in which proteins deposited at splice junctions block m6A over about 100\,nt \citep{uzonyi2023,he2023,luo2023}. Since this mechanism depends on gene structure rather than cell type, it is expected to transfer between cell lines. Predicted RNA folding does not explain the clustering.
% SOURCE: analysis/newdata/annotation_context/ (Ensembl 91, analysis only);
% analysis/newdata/fold_comodification*/.

\paragraph{Why 100\,nt in particular.} It is also the smallest radius that
gives almost every site something to corroborate with: at 50\,nt, 17\% of
sites have no neighbour at all, against 4\% at 100\,nt, where 85\% have two or
more.~(Table~\ref{tab:nbrradius})
% SOURCE: analysis/newdata/distance_bands/neighbours.csv.

\paragraph{Design implication.} Neighbours are limited to about 100--150\,nt.


\paragraph{Radius test, all evaluations.}
\begin{table}[H]
    \centering\small
    \begin{tabular}{lccccc}
        \toprule
        \textbf{Radius} & \textbf{\clone$\to$\cltwo} & \textbf{New genes} & \textbf{Both$\to$\cltwo} & \textbf{Both$\to$\clone} & \textbf{Crossed-test drop} \\
        \midrule
        50\,nt & +0.033 & +0.029 & +0.044 & +0.081 & 0.017 \\
        100\,nt & \textbf{+0.043} & \textbf{+0.054} & \textbf{+0.047} & +0.096 & 0.029 \\
        200\,nt & +0.034 & +0.048 & +0.038 & +0.103 & 0.029 \\
        400\,nt & +0.032 & +0.044 & +0.033 & \textbf{+0.112} & \textbf{0.045} \\
        \bottomrule
    \end{tabular}
    \caption{Radius test (\mdl{h2gcn_aux}, only the radius changed, seed 0).
    Crossed-test drop: loss when the labels change cell line on shared sites, a measure of cell-line specificity. Seed spread $\approx$0.003--0.005.}
    \label{tab:radius}
\end{table}

\paragraph{Which distances a 400\,nt network uses.}
\begin{table}[H]
    \centering\small
    \begin{tabular}{lcccc}
        \toprule
        \textbf{Band removed} & \textbf{\clone-trained: \clone} & \textbf{\clone-trained: \cltwo} & \textbf{Both: \clone} & \textbf{Both: \cltwo} \\
        \midrule
        0--50\,nt & $-0.004$ & $-0.004$ & $-0.009$ & $-0.009$ \\
        50--100\,nt & $-0.005$ & $-0.002$ & $-0.004$ & $-0.004$ \\
        100--200\,nt & $-0.007$ & $-0.004$ & $-0.008$ & $-0.003$ \\
        200--400\,nt & $-0.014$ & +0.001 & $-0.010$ & +0.001 \\
        0--10\,nt & $-0.0001$ & $-0.0010$ & & \\
        10--20\,nt & $-0.0005$ & $-0.0003$ & & \\
        \bottomrule
    \end{tabular}
    \caption{What the trained 400\,nt model uses: change in PR AUC when one
    band of neighbours is removed at prediction time. Distant neighbours help
    only the training line.}
    \label{tab:bands}
\end{table}
% SOURCE: results/band_ablation.csv, band_ablation_v1.csv

\paragraph{Where m6A sits along a gene.}
\begin{figure}[H]
    \centering
    \reportfig[width=0.9\linewidth]{fig_gene_structure.png}
    \caption{Modified-site rate relative to each file's average, by distance
    to the nearest exon junction (left) and position relative to the stop
    codon (right). Ensembl 91; analysis only, never a model input.}
    \label{fig:genestructure}
\end{figure}

\begin{table}[H]
    \centering\small
    \begin{tabular}{lcc}
        \toprule
        \textbf{Context} & \textbf{\clone} & \textbf{\cltwo} \\
        \midrule
        Within 25\,nt of an exon--exon junction & 0.10$\times$ & 0.28$\times$ \\
        25--50\,nt from a junction & 0.13$\times$ & 0.41$\times$ \\
        50--100\,nt & 0.59$\times$ & 0.72$\times$ \\
        100--200\,nt & 2.60$\times$ & 2.17$\times$ \\
        200--400\,nt & 2.87$\times$ & 2.41$\times$ \\
        $>$400\,nt & 1.19$\times$ & 1.25$\times$ \\
        Exon $<$200\,nt & 0.08$\times$ & 0.34$\times$ \\
        Exon 400--800\,nt & 2.46$\times$ & 2.04$\times$ \\
        Internal exon / last exon & 0.25$\times$ / 1.62$\times$ & 0.41$\times$ / 1.65$\times$ \\
        3$'$ UTR / coding sequence & 1.41$\times$ / 0.70$\times$ & 1.51$\times$ / 0.71$\times$ \\
        100\,nt before the stop codon & 1.63$\times$ & 1.54$\times$ \\
        First 100\,nt after the stop codon & 2.68$\times$ & 2.40$\times$ \\
        \bottomrule
    \end{tabular}
    \caption{Positive rate relative to each cell line's average. Validation:
    every labelled 7-mer matches the reference sequence (100\% of a
    19,856-site sample); every annotated stop codon reads TAA/TAG/TGA.}
\end{table}
% SOURCE: analysis/newdata/annotation_context/*.csv

\paragraph{RNA folding test.}
\label{app:folding}

\begin{table}[H]
    \centering\small
    \begin{tabular}{lcccc}
        \toprule
        & \multicolumn{2}{c}{\textbf{Cell line 1}} & \multicolumn{2}{c}{\textbf{Cell line 2}} \\
        \cmidrule(lr){2-3}\cmidrule(lr){4-5}
        \textbf{Distance} & \textbf{LinearPartition} & \textbf{RNAplfold} & \textbf{LinearPartition} & \textbf{RNAplfold} \\
        \midrule
        25--50\,nt & 0.94 [0.72, 1.16] & 1.00 [0.85, 1.17] & 1.15 [0.94, 1.36] & 1.18 [0.99, 1.42] \\
        50--100\,nt & 1.17 [0.94, 1.46] & 0.98 [0.78, 1.18] & 1.35 [1.09, 1.67] & 1.33 [1.08, 1.65] \\
        100--200\,nt & 1.02 [0.67, 1.43] & 0.97 [0.76, 1.22] & 1.03 [0.71, 1.41] & 1.07 [0.80, 1.34] \\
        200--400\,nt & 1.09 [0.56, 1.59] & 0.97 [0.74, 1.18] & 1.19 [0.62, 1.75] & 1.21 [0.99, 1.53] \\
        400--800\,nt & 1.27 [0.53, 2.01] & -- & 1.96 [0.89, 3.32] & -- \\
        \bottomrule
    \end{tabular}
    \caption{Co-modification of fold-linked vs unlinked site pairs at the same
    distance on the same transcripts (1.0 = folding adds nothing; pairing
    probability $\ge$0.1; 6,278 transcripts). RNAplfold \citep{lorenz2011}: window 480, span
    400; LinearPartition \citep{zhang2020linpart}: global. Predicted accessibility of the A shows no
    robust association either (0.88--1.09$\times$).}
\end{table}
% SOURCE: analysis/newdata/fold_comodification/, fold_comodification_linearpartition/

\paragraph{Neighbours within each radius.}
\begin{table}[H]
    \centering\small
    \begin{tabular}{lccccc}
        \toprule
        \textbf{Radius} & \textbf{No neighbours} & \textbf{1} & \textbf{2+} & \textbf{Median} \\
        \midrule
        25\,nt & 42\% & 38\% & 21\% & 1 \\
        50\,nt & 17\% & 29\% & 54\% & 2 \\
        100\,nt & 4\% & 11\% & 85\% & 3 \\
        200\,nt & 1\% & 2\% & 97\% & 7 \\
        400\,nt & 0\% & 0\% & 99\% & 13 \\
        \bottomrule
    \end{tabular}
    \caption{Other candidate sites within each radius (cell line 1; cell line
    2 within 1--2 points). 100\,nt is the smallest radius giving most sites
    several neighbours.}
    \label{tab:nbrradius}
\end{table}

\label{app:capacity-findings}
\paragraph{More Capacity Captures Cell-Line-Specific Patterns}


\paragraph{Motivation.}
The radius test (Section~\ref{sec:reach}) showed that context beyond about 100\,nt helps only the cell line a model was trained on. We asked whether this holds more generally: does giving a model more freedom in how it uses its neighbours help on a cell line it has not seen?

\paragraph{Finding 1: more freedom helps the training cell line, not the unseen one.}
We gave the model more freedom in four ways: letting it learn how much to trust each neighbour (attention; \citealp{velickovic2018}), averaging over the whole transcript, widening the radius to 400\,nt, and a graph transformer \citep{rampasek2022}, which can attend to every site on the transcript. Each improved predictions on the training cell line far more than on the unseen one.~(Table~\ref{tab:capacity}) The graph transformer is the clearest case: it beat H2GCN on cell line 1 (+0.025 PR AUC, $p<0.001$) but not on the unseen cell line ($-0.003$, $p=0.74$). About half of its attention goes to sites more than 150\,nt away, beyond the range where modified sites cluster.~(Figure~\ref{fig:gps})
% SOURCE: results/day_significance.csv; probe_gps.csv, probe_gps_attention.csv;
% the four designs: results/*_nn__xsrc.json (seeds averaged).

\paragraph{Finding 2: additional capacity captures cell-line-specific patterns.}
This does not appear to be ordinary overfitting. When trained and evaluated within cell line 2, the graph transformer again outperforms H2GCN ($p=0.02$); its advantage disappears when training and test cell lines differ (Figure~\ref{fig:ceiling}, right). One possible explanation is transcript-level modification enrichment (Section~\ref{sec:cluster}): the transcripts enriched for m6A differ between cell lines, and whole-transcript attention may learn these differences. In contrast, H2GCN gains little from training on the unseen line's own labels, suggesting that its learned relationships already transfer from cell line 1 (Figure~\ref{fig:ceiling}, right).
% SOURCE: results/data1_arm.csv (p = 0.02; +0.007 p = 0.07; +0.019 p < 0.001).

\paragraph{Finding 3: discordant labels are not resolved by the available reads.}
The two cell lines' labels often disagree: of the sites that appear in both files and are modified in either, 64\% are modified in only one. As a result, even a perfect copy of cell line 1's labels would score only 0.326 PR AUC against cell line 2's, below our model's 0.422 (Figure~\ref{fig:ceiling}, left). At sites where the labels disagree, neither our models nor the raw reads distinguish the two cell lines, so the disagreement lies either in the labelling or below what the reads resolve. Fitting the training cell line more closely cannot recover it.
% SOURCE: results/label_oracle.csv, discordant_sites.csv, discordant_signal.csv.

\begin{figure}[H]
    \centering
    \reportfig[width=0.9\linewidth]{fig_ceiling.png}
    \caption{Left: PR AUC against cell line 2's labels on the 67,320 shared sites, for cell line 1's labels used as predictions (alone, and with ties broken by our model) and for our model. Right: two networks trained on each cell line or both, scored on cell line 2.}
    \label{fig:ceiling}
\end{figure}

\paragraph{Design implication.}
Rather than adding capacity, constrain how the model uses its neighbours, so that it can learn only the parts of the evidence that transfer between cell lines.


The label-copy calculation is a reference predictor on shared sites, not a universal upper bound on achievable performance. Failure of these models and raw-feature checks to distinguish discordant sites does not prove that richer measurements could not do so.


\paragraph{More freedom, more cell-line-specific.}
\begin{table}[H]
    \centering\small
    \begin{tabularx}{\linewidth}{Yccc}
        \toprule
        \textbf{Design (what it adds over H2GCN)} & \textbf{Training line} & \textbf{Unseen line} & \textbf{Cell-line gap} \\
        \midrule
        H2GCN, 50\,nt (reference) & +0.076 & +0.033 & 0.034 \\
        \midrule
        GAT: learned attention over neighbours & +0.055 & +0.002 & 0.052 \\
        Transcript-wide mean (one seed) & +0.093 & +0.019 & 0.070 \\
        400\,nt radius & +0.099 & +0.033 & 0.071 \\
        Graph transformer: attention over the whole transcript & +0.100 & +0.028 & 0.062 \\
        Hand-made neighbour features (LightGBM)$^\dagger$ & +0.065 & +0.016 & -- \\
        \midrule
        \emph{For contrast, constrained designs (Section~\ref{sec:constrain}):} & & & \\
        Scalar messages & +0.075 & +0.055 & 0.021 \\
        Residual + dropout & +0.102 & +0.051 & 0.048 \\
        \bottomrule
    \end{tabularx}
    \caption{PR AUC gain over quantiles\,+\,LightGBM. Training line: trained
    on both cell lines, scored on cell line 1. Unseen line: trained on cell
    line 1 only, scored on cell line 2. Cell-line gap: the gain on cell line
    1 minus the gain on cell line 2, both for models trained on both lines ---
    how much of what a model learns is specific to one line. Three seeds
    averaged unless stated. $^\dagger$Earlier protocol: cross-validation
    within data0 vs scoring on data1's labels (Section~\ref{sec:data}).}
    \label{tab:capacity}
\end{table}
% SOURCE: results/{h2gcn_aux,gat,h2gcn_transcript,h2gcn_aux_r400,gps,scalar_msg,
% res_gate}_nn__xsrc[__seedN].json (cross_source arms; seeds averaged);
% findings.md section 1 (everything).

\paragraph{Graph transformers out of distribution.} ``More capacity hurts out
of distribution'' is not a general law: GPS generalised better than message
passing on 4 of 6 GOOD benchmarks \citep{gui2022,niv2025}, though locality
won on some shifts. Our result is evidence about \emph{label shift between
cell lines} in this domain, and should be phrased that way.

\paragraph{Where the graph transformer attends.}
\begin{figure}[H]
    \centering
    \reportfig[width=0.9\linewidth]{fig_gps_attention.png}
    \caption{Why the graph transformer does not transfer. Left: share of its
    attention by distance from the scored site (about half beyond 150\,nt,
    31\% beyond 400\,nt). Right: the same network rescored with attention cut
    to 150 and 50\,nt; the training line loses 2.6--3$\times$ more (0.569
    $\to$ 0.530 / 0.477 on \clone; 0.403 $\to$ 0.388 / 0.372 on \cltwo). The
    asymmetry, not the size, is the evidence.}
    \label{fig:gps}
\end{figure}

\paragraph{The label ceiling in closed form.}
\label{app:ceiling}
A predictor that is a function of cell line 1's label alone can only rank
$y^{(1)}=1$ sites above $y^{(1)}=0$ sites. On the shared sites,
$P(y^{(2)}{=}1\mid y^{(1)}{=}1)=2136/3291=0.649$ and recall at that cut is
$2136/4853=0.440$, so
$\text{AP}=0.440\times0.649+(1-0.440)\times0.0721=0.326$. This bounds that
class of predictors exactly; it is not a bound on every model.

\begin{table}[H]
    \centering\small
    \begin{tabular}{lcc}
        \toprule
        \textbf{Score for cell line 2's labels (shared sites, 7.2\% positive)} & \textbf{PR AUC} & \textbf{ROC AUC} \\
        \midrule
        Cell line 1's true labels & 0.326 & 0.711 \\
        Cell line 1's true labels, ties broken by \mdl{h2gcn_aux} & 0.485 & 0.837 \\
        \mdl{h2gcn_aux} trained on cell line 1 (seeds averaged) & 0.422 & 0.823 \\
        \bottomrule
    \end{tabular}
    \caption{Label oracle. The reverse direction: cell line 2's labels score
    0.30 on cell line 1's (0.63 tie-broken), against 0.59 for our model.}
\end{table}

\begin{table}[H]
    \centering\small
    \begin{tabular}{lrr}
        \toprule
        \textbf{Shared sites (67,320)} & \textbf{Count} \\
        \midrule
        Modified in both & 2,136 \\
        Modified in cell line 1 only & 1,155 \\
        Modified in cell line 2 only & 2,717 \\
        Unmodified in both & 61,312 \\
        \bottomrule
    \end{tabular}
    \quad
    \begin{tabular}{lcc}
        \toprule
        \textbf{Model (seed 0)} & \textbf{Which-line AUC} & \textbf{Rank corr.} \\
        \midrule
        Quantiles + LightGBM & 0.52 & 0.72 \\
        DeepSet & 0.50 & 0.88 \\
        \mdl{h2gcn_aux} & 0.48 & 0.91 \\
        \mdl{res_gate} & 0.48 & 0.91 \\
        \mdl{scalar_msg} & 0.48 & 0.91 \\
        \mdl{gps} & 0.45 & 0.93 \\
        \bottomrule
    \end{tabular}
    \caption{Left: label agreement on shared sites. Right: at the 3,872
    discordant sites, can a model tell which cell line carries the
    modification from its reads? (0.5 = no.) Rank corr.: a site's rank in
    the two files' reads.}
\end{table}

\begin{table}[H]
    \centering\small
    \begin{tabular}{lccc}
        \toprule
        \textbf{Linear direction} & \textbf{Within \clone} & \textbf{Within \cltwo} & \textbf{Which-line, discordant} \\
        \midrule
        Fitted on cell line 1 & 0.818 & 0.712 & 0.519 \\
        Fitted on cell line 2 & 0.819 & 0.714 & 0.516 \\
        None: centre-position current & 0.301 & 0.385 & 0.542 \\
        \bottomrule
    \end{tabular}
    \caption{Model-free raw-signal check (AUC): per-site means and spreads of
    the 9 read features, centred per 7-mer, projected on a direction fitted on
    non-shared sites.}
\end{table}
% SOURCE: results/label_oracle.csv, discordant_sites.csv, discordant_signal.csv

\phantomsection
\subsection*{Constraint ablations: supporting findings\texorpdfstring{\appmark{app:constraint-findings}}{}}
\label{app:constraint-findings}
\paragraph{Constraining the Neighbours}


\paragraph{Motivation.}
Section~\ref{sec:capacity} concluded that the model should be constrained rather than given more capacity. Here we ask how: which constraints on the way neighbours are used make the graph transfer better to an unseen cell line?

We compared distance weighting, residual corrections with neighbour dropout,
and scalar messages, using the definitions in Methods and the controls in
Supplementary Methods.\applink{app:designs}

\begin{table}[H]
    \centering\small
    \begin{tabular}{lcc}
        \toprule
        \textbf{Design} & \textbf{Unseen line} & \textbf{Training line} \\
        \midrule
        H2GCN, free neighbour mean within 50\,nt (Section~\ref{sec:graph}) & +0.033 & +0.076 \\
        Distance weighting & +0.043 & +0.096 \\
        \textbf{Residual + dropout} & \textbf{+0.051} & \textbf{+0.102} \\
        \textbf{Scalar messages + dropout} & \textbf{+0.054} & +0.088 \\
        Graph transformer, for contrast (Section~\ref{sec:capacity}) & +0.028 & +0.100 \\
        \bottomrule
    \end{tabular}
    \caption{PR AUC gain over quantiles\,+\,LightGBM, mean of three seeds. Unseen line: trained on cell line 1, scored on cell line 2; training line: trained on both, scored on cell line 1.~(Table~\ref{tab:batch})}
    \label{tab:batch_short}
\end{table}
% SOURCE: analysis/representation/results/day_summary.md (3 seeds).

\paragraph{Finding 1: constraining transfers better.}
Every constrained design beats the free neighbour mean on the unseen cell line, by +0.010 to +0.022 PR AUC: residual + dropout by +0.017 and scalar messages by +0.020 (each $p<0.001$).~(Table~\ref{tab:significance}) Since the unseen line is always the weaker of the two, the selection measure of Section~\ref{sec:data} rises from +0.033 to +0.051--0.055. Simply reaching further does not explain this: a plain 100\,nt radius gains only +0.004 over 50\,nt, against +0.010 for distance weighting.
% SOURCE: results/day_significance.csv; day_summary.md (matched differences).

\paragraph{Finding 2: neighbour dropout accounts for most of the residual design's advantage.}
Removing dropout reduces the residual design's gain on the unseen line, whereas removing its weighting of neighbours by read count has no measurable effect.~(Table~\ref{tab:batch}) Dropout prevents the model from relying on any single neighbour. Combined with the residual form, it also keeps predictions stable when neighbours are missing, as in partial transcripts: a site with no neighbours keeps its own-reads score.
% SOURCE: day_summary.md (res_nodrop, res_nogate, 3 seeds); probe_gate_params.csv.

\paragraph{Finding 3: distance weighting helps, but not through its measured profile.}
Reversing the weights of the three distance bands keeps most of the gain.~(Table~\ref{tab:batch}) The benefit appears to come from reaching about 150\,nt while treating near and far neighbours separately, rather than from the precise shape of the co-modification curve.
% SOURCE: day_summary.md (fk_band vs fk_band_shuffled, 3 seeds).

\paragraph{Design implication.}
The final model carries the two strongest constrained designs, residual + dropout and scalar messages + dropout (Section~\ref{sec:model}).


\paragraph{Constrained designs, all evaluations.}
\begin{table}[H]
    \centering\footnotesize\setlength{\tabcolsep}{4pt}
    \begin{tabularx}{\linewidth}{lccccY}
        \toprule
        \textbf{Design} & \textbf{Unseen line} & \textbf{New genes} & \textbf{Both$\to$\cltwo} & \textbf{Both$\to$\clone} & \textbf{What it tests} \\
        \midrule
        H2GCN + own-reads head & +0.033 & +0.026 & +0.042 & +0.076 & Reference: free mean within 50\,nt \\
        \quad hard 100\,nt radius & +0.037 & +0.045 & +0.043 & +0.087 & Reach alone \\
        Distance bands & +0.045 & +0.053 & +0.045 & +0.091 & Bands, weights from labels \\
        \quad control: reversed weights & +0.043 & +0.055 & +0.044 & +0.092 & Do the measured weights matter? \\
        Distance kernel & +0.043 & +0.045 & +0.052 & +0.096 & Smooth $e^{-d/\lambda}$ \\
        \textbf{Residual + dropout} & \textbf{+0.051} & +0.053 & +0.053 & \textbf{+0.102} & Own score + gated correction \\
        \quad control: no gate & +0.049 & +0.045 & +0.059 & +0.108 & Is the depth gate needed? \\
        \quad control: no dropout & +0.040 & +0.040 & +0.050 & +0.102 & Is dropout needed? \\
        Scalar messages & \textbf{+0.055} & +0.052 & +0.054 & +0.075 & Neighbours pass one number \\
        \quad control: random neighbours & +0.017 & +0.016 & +0.028 & +0.057 & Is it corroboration? \\
        \textbf{Scalar messages + dropout} & +0.054 & +0.050 & \textbf{+0.059} & +0.088 & Both ingredients \\
        Graph transformer & +0.028 & +0.021 & +0.038 & +0.100 & Capacity (Section~\ref{sec:capacity}) \\
        \bottomrule
    \end{tabularx}
    \caption{Every constrained design and control behind Table~\ref{tab:batch_short}. PR AUC gain over quantiles\,+\,LightGBM, mean of 3 seeds (seed SD
    0.001--0.005 for most). ``Unseen line'': trained on cell line 1, scored
    on cell line 2. ``New genes'': data1 genes absent from data0.
    Indented rows are controls.}
    \label{tab:batch}
\end{table}
% SOURCE: analysis/representation/results/day_summary.md; 3-seed controls from findings 12F.

\paragraph{Paired comparisons with H2GCN.}
\begin{table}[H]
    \centering\footnotesize\setlength{\tabcolsep}{3pt}
    \resizebox{\linewidth}{!}{%
    \begin{tabular}{lcccc}
        \toprule
        \textbf{vs \mdl{h2gcn_aux}} & \textbf{Unseen line} & \textbf{New genes} & \textbf{Both$\to$\cltwo} & \textbf{Both$\to$\clone} \\
        \midrule
        \mdl{res_gate} & +0.017 [.011, .023] & +0.028 [.012, .042] & +0.009 [.003, .016] & +0.019 [.010, .028] \\
        \mdl{scalar_msg} & +0.020 [.013, .028] & +0.024 [.007, .041] & +0.010 [.002, .017] & $-0.007$ [$-.017$, .004] \\
        \mdl{res_gate} + \mdl{scalar_msg} & +0.024 [.018, .029] & +0.029 [.012, .045] & +0.016 [.009, .022] & +0.019 [.010, .027] \\
        \mdl{h2gcn_aux_r100} & +0.005 [$-.002$, .011] & +0.022 [.004, .039] & $-0.000$ [$-.008$, .006] & +0.009 [$-.001$, .018] \\
        \mdl{gps} & $-0.003$ [$-.011$, .005] & $-0.003$ [$-.018$, .012] & $-0.002$ [$-.009$, .005] & +0.025 [.015, .035] \\
        \mdl{scalar_random} (seed 0) & $-0.018$ [$-.026$, $-.010$] & $-0.013$ [$-.032$, .007] & $-0.016$ [$-.024$, $-.008$] & $-0.024$ [$-.036$, $-.013$] \\
        \bottomrule
    \end{tabular}}
    \caption{Paired PR AUC difference from \mdl{h2gcn_aux} on identical sites,
    seeds rank-averaged, 95\% interval from 1,000 gene resamples.}
    \label{tab:significance}
\end{table}
% SOURCE: results/day_significance.csv

\phantomsection
\subsection*{Robustness: supporting findings\texorpdfstring{\appmark{app:robust-findings}}{}}
\label{app:robust-findings}
\paragraph{Robustness}


\paragraph{New sequencing runs.} Two further SG-NEx HCT116 runs --- the same
RNA sample re-sequenced, and a separately grown batch of cells --- carry
data0's labels at all 121,838 sites. Scored with the fold model that held
out each gene, using only the new run's reads, a close variant of the final
model (scalar messages without dropout in place of with) loses 0.009--0.015
PR AUC against data0's own run, and single networks up to
0.025.~(Table~\ref{tab:hct116runs}) This is larger than the seed spread
(0.002--0.005), so a new run does cost a little, but it is far smaller than
the cross-cell-line drop. The biological replicate scores no worse than
the technical one.

\paragraph{Low read depth.} Every training site has $\ge$20 reads, so lower
depths were simulated by thinning each held-out site's reads, and its
neighbours', at random (Table~\ref{tab:depth}; Figure~\ref{fig:depth}). Each
read is a separate molecule, so this is a fair simulation of shallower
sequencing, though real low-depth sites come from lowly expressed genes,
which may differ in other ways. At 1--3 reads, the depths most SG-NEx sites
have, the final model keeps about twice LightGBM's PR AUC. Most of this
advantage comes from the read encoder (Section~\ref{sec:reads}, Finding 3):
on cell line 2 at 1 read, the network without neighbours already scores
0.225 against LightGBM's 0.136, and equally thinned neighbours add only
0.018. Corroboration adds most when a site's neighbours are better covered
than the site itself: with neighbours at full depth, the same 1-read site
scores 0.367.~(Table~\ref{tab:swap})


% SOURCE: analysis/representation/results/depth_ensembles.csv

\paragraph{Known modification fractions (data2).} One synthetic RNA was
sequenced as seven samples with a known share of modified molecules (0, 25,
50, 70, 75, 95, 100\%) at the same 189 DRACH positions; the model never
trained on it. Scores separate unmodified from modified RNA sharply (mean
submitted score 0.15 at 0\%, 0.44 at 25\%, 0.57 at 50\%) and then plateau.~(Figure~\ref{fig:data2}) Here every neighbour shares the
site's modification level by construction, and adding the neighbour step
raises the ROC AUC for telling 0\% from any modified sample from 0.75 to 0.87
(residual) and 0.77 to 0.90 (scalar messages). This graph is perfectly
homophilous, unlike real transcripts (Section~\ref{sec:graph}), so it is a
best case for corroboration rather than a typical one. The score detects the \emph{presence} of m6A; it does
not measure the modified fraction (Section~\ref{sec:limits}).
% SOURCE: analysis/newdata/data2_scores.py (shipped model); means of score_rank
% per fraction from .cache/representation/data2_scores.csv (course data, not
% committed); medians in analysis/representation/results/data2_score_summary.csv.

% =============================================================================

\paragraph{Further HCT116 sequencing runs.}
\begin{table}[H]
    \centering\small
    \begin{tabular}{lccc}
        \toprule
        \textbf{Model (trained on)} & \textbf{Rep.\ 3 run 4} & \textbf{Rep.\ 4 run 3} & \textbf{data0's run} \\
        \midrule
        \mdl{h2gcn_aux} (\clone) & 0.523 & 0.533 & 0.544 \\
        \mdl{res_gate} (\clone) & 0.562 & 0.567 & 0.584 \\
        \mdl{scalar_msg} (\clone) & 0.539 & 0.549 & 0.560 \\
        \mdl{gps} (\clone) & 0.554 & 0.551 & 0.569 \\
        \mdl{h2gcn_twohead_aux} (both) & 0.561 & 0.560 & 0.586 \\
        Intermediate ensemble (both, seed 0) & 0.575 & 0.578 & 0.595 \\
        Two-head + \mdl{res_gate} + \mdl{scalar_msg} (both) & 0.584 & 0.590 & 0.599 \\
        \bottomrule
    \end{tabular}
    \caption{PR AUC on two further SG-NEx HCT116 runs with data0's labels,
    each site scored by the fold model that held out its gene. Every site has
    $\ge$20 reads in both runs (median 42), so this tests run-to-run
    variation, not depth. The last row differs from the final model only in
    using scalar messages without dropout.}
    \label{tab:hct116runs}
\end{table}
% SOURCE: analysis/representation/results/hct116_runs.csv

\paragraph{PR AUC against read depth.}
\begin{figure}[H]
    \centering
    \reportfig[width=0.85\linewidth]{fig_depth.png}
    \caption{PR AUC against reads per site, simulated by thinning each site's
    and its neighbours' reads (Table~\ref{tab:depth}). Grey band: 1--3 reads,
    typical for SG-NEx. The intermediate model is within 0.014 of the final
    model at every depth.}
    \label{fig:depth}
\end{figure}
% SOURCE: analysis/representation/results/depth_ensembles.csv

\paragraph{Known modification fractions (data2).}
\begin{figure}[H]
    \centering
    \reportfig[width=0.95\linewidth]{fig_data2_fractions.png}
    \caption{In-vitro dose series (data2): 189 positions, seven samples with
    known modified fractions; the same shipped constrained networks with
    (left) and without (right) their neighbour step. The submitted score
    (mean per sample) rises from 0.15 at 0\% to 0.44 at 25\% and 0.57 at
    50\%, then plateaus at 0.56--0.60. Every neighbour shares the
    site's level by construction here, a best case for corroboration; sites
    have 600--1,200 reads, outside the training range. The full ensemble's
    Spearman with fraction is 0.41.}
    \label{fig:data2}
\end{figure}

\phantomsection
\subsection*{GLORI agreement by depth, and the m6Anet comparison\texorpdfstring{\appmark{app:glori-results}}{}}
\label{app:glori-results}
\label{app:glori-depth-results}
PR AUC is average precision. Lift is PR AUC divided by the observed GLORI
positive rate. The motif-only baseline assigns each DRACH 5-mer its mean
modification rate across the two labelled training lines. Top-1\% overlap
is measured within each line and depth band. These comparisons use the same
inferred labels for both rankings. Pooled comparisons can mix coverage effects
with within-band ranking quality.

\begin{table}[H]
\centering\footnotesize\setlength{\tabcolsep}{3pt}
\begin{tabular}{llrrrrr}
\toprule
Line & Pooled reads & Sites & Positive rate & Model AP & Motif AP & Top 1\% overlap \\
\midrule
A549 & 1-2 & 140,503 & 3.3\% & 0.135 & 0.073 & 19.2\% \\
A549 & 3-9 & 171,655 & 4.3\% & 0.212 & 0.097 & 25.9\% \\
A549 & 10-19 & 83,513 & 5.0\% & 0.264 & 0.113 & 29.0\% \\
A549 & 20+ & 122,873 & 6.1\% & 0.366 & 0.150 & 45.4\% \\
H9 & 1-2 & 104,168 & 1.0\% & 0.041 & 0.021 & 5.4\% \\
H9 & 3-9 & 145,018 & 2.2\% & 0.103 & 0.044 & 11.0\% \\
H9 & 10-19 & 97,075 & 3.0\% & 0.149 & 0.063 & 12.8\% \\
H9 & 20+ & 224,439 & 5.1\% & 0.300 & 0.107 & 39.5\% \\
HepG2 & 1-2 & 21,904 & 0.4\% & 0.022 & 0.012 & 2.7\% \\
HepG2 & 3-9 & 23,740 & 1.2\% & 0.086 & 0.035 & 14.3\% \\
HepG2 & 10-19 & 18,377 & 1.9\% & 0.138 & 0.057 & 15.2\% \\
HepG2 & 20+ & 98,743 & 4.6\% & 0.381 & 0.141 & 51.0\% \\
HCT116* & 1-2 & 19,490 & 0.2\% & 0.015 & 0.008 & 3.6\% \\
HCT116* & 3-9 & 18,899 & 1.0\% & 0.072 & 0.025 & 8.5\% \\
HCT116* & 10-19 & 14,207 & 1.6\% & 0.097 & 0.043 & 12.7\% \\
HCT116* & 20+ & 46,980 & 2.8\% & 0.194 & 0.085 & 21.5\% \\
\bottomrule
\end{tabular}
\caption{GLORI agreement by pooled read depth. HCT116* excludes training sites but is not an unseen cell line. Top-1\% overlap is for our model.}
\label{tab:glori-depth}
\end{table}

\label{app:glori-comparison}
\paragraph{GLORI-based m6Anet comparison.}
We used m6Anet 2.1.0 with pretrained \texttt{HCT116\_RNA002}. Both models
were evaluated on matched genomic sites in A549, H9 and HepG2, unseen by
both. m6Anet's input was capped at 50 randomly sampled reads per site (seed
4262); our model used all pooled reads. Depth bands use uncapped read counts.
Below 20 reads, the inference minimum was lowered to one. This extended
setting should not be confused with m6Anet's default coverage. Differences
and central 95\% intervals use 200 resamples of whole genes.

Under the default depth requirement, eligible sites comprise 23.7\% of the
graded A549 candidates, 39.3\% of H9 and 60.7\% of HepG2. These contain
31.9\%, 61.5\% and 86.2\% of the matched GLORI positives, respectively.
The extended run produced scores for all but 34 A549 and 60 H9 candidates,
and all HepG2 candidates. All-depth PR AUC differences describe a pooled
ranking, not a matched-depth advantage or proof of a particular transfer
mechanism. A549 top-1\% GLORI overlap at 20+ reads favoured m6Anet
(49.1\% versus 45.4\%), whereas H9 and HepG2 favoured our model
(39.5\% versus 32.0\%; 51.0\% versus 33.7\%).

\begin{table}[H]
\centering\footnotesize\setlength{\tabcolsep}{3pt}
\begin{tabular}{llrrl}
\toprule
Line & Pooled reads & Our AP & m6Anet AP & Difference [95\% interval] \\
\midrule
A549 & 1-2 & 0.135 & 0.114 & +0.021 (+0.016 to +0.026) \\
A549 & 3-9 & 0.212 & 0.205 & +0.008 (+0.003 to +0.014) \\
A549 & 10-19 & 0.264 & 0.275 & -0.011 (-0.018 to -0.004) \\
A549 & 20+ & 0.366 & 0.357 & +0.008 (-0.001 to +0.017) \\
A549 & all depths & 0.270 & 0.201 & +0.069 (+0.064 to +0.074) \\
H9 & 1-2 & 0.041 & 0.037 & +0.004 (+0.001 to +0.007) \\
H9 & 3-9 & 0.103 & 0.103 & +0.000 (-0.004 to +0.005) \\
H9 & 10-19 & 0.149 & 0.158 & -0.008 (-0.013 to -0.004) \\
H9 & 20+ & 0.300 & 0.262 & +0.038 (+0.032 to +0.044) \\
H9 & all depths & 0.233 & 0.146 & +0.088 (+0.082 to +0.094) \\
HepG2 & 1-2 & 0.022 & 0.016 & +0.006 (-0.002 to +0.033) \\
HepG2 & 3-9 & 0.086 & 0.065 & +0.022 (+0.008 to +0.038) \\
HepG2 & 10-19 & 0.138 & 0.126 & +0.012 (-0.001 to +0.028) \\
HepG2 & 20+ & 0.381 & 0.278 & +0.103 (+0.087 to +0.119) \\
HepG2 & all depths & 0.341 & 0.173 & +0.169 (+0.153 to +0.185) \\
\bottomrule
\end{tabular}
\caption{Matched-site GLORI comparison. Below 20 reads, m6Anet uses the modified inference minimum. Differences reflect unrounded values.}
\label{tab:glori-m6anet-depth}
\end{table}

\phantomsection
\subsection*{GLORI fractions and cross-assay agreement\texorpdfstring{\appmark{app:glori-fractions}}{}}
\label{app:glori-fractions}
Among GLORI-positive sites with 20+ pooled reads, score--fraction Spearman
correlations were 0.585, 0.546 and 0.639 for our model and 0.623, 0.606 and
0.708 for m6Anet (A549, H9, HepG2). For our model within GGACT, corresponding
correlations were 0.493, 0.346 and 0.642, based on 1,186, 1,320 and 1,073
positive sites. HCT116 correlations were 0.677 overall at 20+ reads and
0.674 within GGACT (321 sites), after training-site exclusion. Restriction
to GLORI positives and coverage selection limits generalisation of these
correlations; association is not calibration or validation of absolute
modification-fraction estimates.

For assay agreement, 48,620 HCT116 training sites in GLORI-covered genes
were compared separately: course-label precision against GLORI was 0.477,
recall 0.443, and GLORI positive rate 6.6\%. This is a comparison of labels,
not a held-out model evaluation. It cannot be used as a performance ceiling
or directly compared with precision in other cell lines and candidate sets.

\paragraph{Provenance.}
The findings above are recorded in \texttt{report/sgnex\_glori\_findings.md}.
Mapping and motif comparisons use \texttt{analysis/sgnex/glori/}; matched
m6Anet comparisons use \texttt{analysis/sgnex/m6anet\_glori/}; gene-group
analyses use \texttt{analysis/sgnex/gene\_classes/}. Each directory contains
its generated summary and metrics or effects table. The corresponding
scripts are \texttt{glori\_check.py}, \texttt{m6anet\_glori\_compare.py} and
\texttt{gene\_classes.py} under \texttt{analysis/sgnex/}.
% Samples processed: 23 across seven cell lines (analysis/sgnex/label_mapping_check.csv
% and the scoring logs agree; analysis/sgnex/README.md's "22" is stale). The text
% reports by cell line, not by sample count.

\phantomsection
\subsection*{Why scores plateau: presence, not fraction\texorpdfstring{\appmark{app:presence}}{}}
\label{app:presence}
Trained with cross-entropy on binary labels, the ideal output is
$s^*(x)=P(Y=1\mid x)$, the probability that the site is modified at all.
A site with modified fraction $f$ yields reads from $(1-f)P_0+fP_1$; the
evidence for ``modified at all'' grows as $n\cdot D(f)$ with
$D(f)\approx\tfrac{f^2}{2}\chi^2(P_1\|P_0)$ for small $f$. With 600--1,200
reads, $f=0.25$ already gives overwhelming evidence and the score
saturates (Figure~\ref{fig:data2}).

\phantomsection
\subsection*{Missing neighbours and partial transcripts\texorpdfstring{\appmark{app:partial-transcripts}}{}}
\label{app:partial-transcripts}
 A site with no
    neighbour within range relies on its own reads alone. With every
    neighbour removed at prediction time, the intermediate ensemble falls to
    0.391 PR AUC on cell line 2, about LightGBM's level. The two-head network,
    which leans on its transcript channel, pulls it down (0.327 alone, against
    0.405 for H2GCN alone). The test
    data may contain partial transcripts. Neighbour dropout in the
    constrained designs trains for this case, but a dedicated safeguard was
    designed and not shipped.

\end{document}
